
\section{Middle primes II: reweighting the principal law on seven intervals}
\label{sec:positive-roots}

\paragraph{What this section proves and how to read it.}
On seven middle intervals from $2^{30}$ onward, the uniform
twelve-label law has spectral bound $g\to\sqrt{30}$, which forces
completion weights too large for $\Theta<1$. This section
reweights it by an explicit nonnegative function (a \emph{tilted
law}, called a \emph{positive root} in the text) that preserves
support and one-slot marginals and improves the bound to
$g(\lambda)<\sqrt{30}$ (where $\lambda=2\ulaw/(1+\ulaw^2)$ in
terms of the tilt parameter $\ulaw$ of
\cref{def:positive-root}, and $g(\lambda)$ is the leading
coefficient defined there), at the price of a larger bound at a
second
cut and a changed one-half marginal, both bounded here. The other
$26$ intervals use the uniform law of \cref{sec:middle-reference};
the choice of which intervals are reweighted is part of the
construction (\cref{sec:design}). Inputs: \cref{sec:middle-reference}
and Part~IV trace estimates. Outputs: the two cut bounds and
marginal controls used in \cref{sec:middle-frame}. A first reading
should follow the definition and the two bound statements; the
pair-score calculations are reference material.
On seven pieces of the middle range, from $2^{30}$ onward, the two-layer
root is a positive reweighting of the uniform twelve-label law. Six pieces
use individual squares of the three quartic pair scores; the last uses
39 coherent traces, with a positive floor. The support, invariance under the chamber
reflections and layer swap of \cref{lem:layer12-geometry}(a) and uniform
one-slot marginals are preserved. We bound the
chamber and layer operators in the fixed uniform induced-half charts
used by \cref{sec:middle-frame}, and give the actual marginal variance
and range needed by \cref{lem:record:root-marginal}(iv).
The pair-score toolkit is proved first; then come the sharper
collision bound on the two protected subspaces, the interpolation
between the pair bounds for the actual (physical) kernel
(\cref{imp:env:physical-interpolation}), the
layer remainder summed without expanding all monomials
(\cref{imp:env:layer-resummation}), and the
finite bounds for the coherent law. The seven endpoint caps
$g,c^+$ are in \cref{cert:positive-root-caps}. These caps
determine the completion
weights $w_{\rm D}$ and $v_{\rm L}$.

The reason for the tilt is visible in the leading matrices. The negative cap of L20 of
the uniform root tends to $\sqrt{30}$, the norm of the complete star matrix of its soft
kernel (\cref{lem:l20-star-matrix}), and through the completion inequality
$w_{\rm D}\ge d$ of \cref{lem:front12:reference}(iii), where $d$ is the block cap
of L20, this cap fixes the first-order weights $w_{\rm D}=w_{\rm S}$ of the
reference. The linear part $2\ulaw\Hbar$ of a tilt adds to that matrix a coupling of the
pairs of opposite base labels, and the eigenvector of the least eigenvalue has a component along
this coupling; the least eigenvalue rises to $-g(\lambda)$ with $g(\lambda)<\sqrt{30}$
(\cref{lem:root-pair-matrix}). This gain is at order one in the coefficient of $r$. The
costs are of order $r$: the layer cut of the tilted root has a larger cap,
which enters the reference only through the larger one-layer weight $v_{\rm L}$
it forces (\cref{rem:middle-weights}), and the half marginal of the root changes,
which the restorations pay through its range ratio and its chi-square.

\emph{Conventions.} Those of \cref{sec:middle-reference}. Throughout $p\equiv3\pmod4$,
$q=1/p$, $r=p^{-1/2}$, every space carries a probability measure, $\ulaw\ge0$ is
rational, and $P_0$ is the projection onto mean-zero functions.
In \cref{imp:local:collision-gap} the same centering $P_0$ is written $Q$,
with protected subspaces $P_{\rm c},P_{\rm p}$ defined there. For an operator $T$ with
a kernel on a product of two spaces and a function $f$ on that product, $T\circ f$ is the
operator whose kernel is the pointwise product of the kernel of $T$ with $f$, as in
\cref{lem:swapped-rectangle,lem:tilt-marginals}.

\subsection{The tilted law \texorpdfstring{$\Gamma_{p,\ulaw}$}{Gamma\_(p,u)}}\label{subsec:positive-root-law}

On seven middle intervals the uniform twelve-label law is
reweighted (tilted) by an explicit sum of squares, hence by a
nonnegative weight with the same support. This one-parameter
family is the positive root; the definition fixes the weight and
the parameter value used on each interval.

\begin{definition}[Positive roots]\label{def:positive-root}
Let $x$ be a labelling of the twelve vertices $(u,\ell)$ of $G_{\rm L,2}$ at the
prime $p$. For $j\in\{0,1,2\}$ let
\[
 \Hs j(x)=r\sum_{z\in\Fp}\chi\Bigl(\prod_{u\in\{j,j+3\}}\ \prod_{\ell\in\{0,1\}}
 (z-x_{u,\ell})\Bigr)
\]
be the normalized trace score of the opposite base pair $\{j,j+3\}$, and
$\Hbar=(\Hs0+\Hs1+\Hs2)/3$. For rational $\ulaw\ge0$ the
\emph{individual-square positive root} is
\begin{equation}\label{eq:positive-root}
 \Gamma_{p,\ulaw}=\frac{w_\ulaw}{\E w_\ulaw}\,\lambda_p(G_{\rm L,2}),\qquad
 w_\ulaw=\frac13\sum_{j=0}^2(1+\ulaw\Hs j)^2,
\end{equation}
the expectation being under $\lambda_p(G_{\rm L,2})=\Gamma_{p,0}$. Thus
\[
 w_\ulaw=1+2\ulaw\Hbar+\ulaw^2 V,
 \qquad V=\frac13\sum_j\Hs j^2.
\]
We put $\gamma=\ulaw^2$, $s=1+\gamma$ and $\lambda=2\ulaw/s$,
the \emph{strength} of the tilt. On a block,
$\Gamma_{\Blk,\ulaw}=\bigotimes_{p\in\Blk}\Gamma_{p,\ulaw}$. The selected construction uses individual squares on the six successive
middle pieces starting at $2^{30},2^{31},2^{32},2^{34},2^{36},2^{40}$, with parameters
$711/4096$, $1021/4096$, $1455/4096$, $833/1024$, $3851/4096$, $4049/4096$,
respectively. On $[2^{45},P_3)$ it uses instead
\[
 \Gamma_{p,\ulaw}=
 \frac{1/100+(1+\ulaw H_{\mathcal D_{39}})^2}
 {\E[1/100+(1+\ulaw H_{\mathcal D_{39}})^2]}
 \,\lambda_p(G_{\rm L,2}),\qquad \ulaw=67/40,
\]
where $H_{\mathcal D_{39}}=\sum_{S\in\mathcal D_{39}}c_SH_S$
with coefficients $c_S\in\{\pm1/15\}$ and $H_S$ the trace
function of the support $S$; the list $\mathcal D_{39}$ of $39$
subsets (``supports'') of the twelve label positions, with the
signs $c_S$, is defined in \cref{imp:local:filtered-coherent}. All expectations normalising
these weights are under the uniform law $\lambda_p(G_{\rm L,2})$.
Below $2^{30}$ the principal law remains uniform.
\end{definition}

Why squares of the individual scores, rather than the square
$(1+\ulaw\Hbar)^2$ of their average? The two weights have the same
linear term $2\ulaw\Hbar$, but the square of the average also
contains the mixed products $\Hs j\Hs k$ ($j\ne k$), whose cut
operators are harder to bound; the individual-square weight
avoids them. On the last interval such products are instead
included in a controlled form: the score $H_{\mathcal D_{39}}$ is
one symmetrised list of trace functions, invariant under all
chamber reflections and the layer swap, so that the same tilted
law can be bounded at every sorting cut. It is called the
\emph{coherent} law for this reason, and the list of supports
with their coefficients is called its \emph{dictionary}.

The three base pairs are $\{0,3\}$, $\{1,4\}$ and $\{2,5\}$: each consists of a base of the
even component $\{0,2,4\}$ of $G_{\rm L,2}$ and a base of the odd component $\{1,3,5\}$, and
a score is a quartic character sum over the four labels of the pair, two in each layer.

\begin{lemma}[Properties of the positive roots]\label{lem:positive-root-properties}
Let $\Gamma_{p,\ulaw}$ be the individual-square root of
\eqref{eq:positive-root}. The coherent root has the separate admissibility statement of
\cref{lem:imp:coherent39-law}.
\begin{enumerate}[label=(\alph*)]
\item On the support of $\Gamma_{p,0}$ every score lies in $r\Z$ and satisfies
 $|\Hs j|\le2+r$, and $w_\ulaw>0$. Hence $\Gamma_{p,\ulaw}$ is a probability law with the
 support of $\Gamma_{p,0}$, and it accepts the six queries.
\item The chamber reflections $u\mapsto k-u$ ($k\in\{1,3,5\}$) and the translation
 $u\mapsto u+2$ permute the three scores, and the layer swap $\ell\mapsto1-\ell$ and every
 map $z\mapsto\alpha z+\beta$ of the labels with $\alpha\ne0$ fix each of them. Hence
 $\Gamma_{p,\ulaw}$, and $\Gamma_{\Blk,\ulaw}$ on a block, is invariant under the maps of
 \cref{lem:layer12-geometry}(a) and under the square-affine maps, and every slot is
 uniform.
\item At the chamber cut $\vartheta_0$, with plus half $(a,b,c)$ on the objects
 $0,1,2$ and reflected half $(a',b',c')$ on the objects $5,4,3$, the scores pair $a$
 with $c'$, $b$ with $b'$ and $c$ with $a'$; the reflection exchanges $\Hs0$ and $\Hs2$
 and fixes $\Hs1$. At the layer cut $\iota$ each score selects in each layer one label of
 each of the two components of $G_{\rm L,2}$, at the same bases.
\end{enumerate}
The half marginals of $\Gamma_{p,\ulaw}$ differ from the uniform induced half laws at
order $r$, as quantified in \cref{prop:root-chamber,prop:root-layer}; all caps
below are proved in those fixed uniform half charts, which are the
diagonal measures of the completion.
\end{lemma}

\begin{proof}
(a) On the support all twelve labels are distinct, so the four labels of a base pair are
distinct. The sum in $\Hs j$ is an integer, and \cref{thm:weil}(b) with $t=4$ bounds its
modulus by $2\sqrt p+1$; hence $\Hs j\in r\Z$ and $|\Hs j|\le2+r$. Let $\ulaw>0$. If
$1+\ulaw\Hs j=0$ with $\Hs j=rn$, $n\in\Z$, then $\sqrt p=-\ulaw n$ would be rational;
so every square in $w_\ulaw$ is strictly positive on the support (and
$w_0=1$), $\E w_\ulaw>0$, and $\Gamma_{p,\ulaw}$ has the support of $\Gamma_{p,0}$.
Acceptance of a query is a property of the support, and $\Gamma_{p,0}$ accepts the six
queries.

(b) $\Hs j$ depends only on the four labels on the vertices $(u,\ell)$ with
$u\in\{j,j+3\}$, symmetrically. The map $u\mapsto k-u$ sends $\{j,j+3\}$ to
$\{k-j,k-j+3\}$ and $u\mapsto u+2$ sends it to $\{j+2,j+5\}$ (indices mod $6$), both base
pairs; the layer swap fixes the vertex set of each base pair. Under $z\mapsto\alpha z+\beta$
the substitution $z=\alpha z'+\beta$ in the defining sum multiplies the product by
$\alpha^4$, and $\chi(\alpha^4)=1$. The weight $w_\ulaw$ is a symmetric function of
$(\Hs0,\Hs1,\Hs2)$, so every map that preserves $\Gamma_{p,0}$ and permutes the scores
preserves $w_\ulaw$ and $\Gamma_{p,\ulaw}$. The maps of \cref{lem:layer12-geometry}(a) are
automorphisms of $G_{\rm L,2}$ and preserve its uniform labelling law, and so do the
square-affine maps, which preserve strict squareness of differences; a block law is a
product. Translations act transitively on the labels, so every slot is uniform
(\cref{lem:marginals}).

(c) At $\vartheta_0$ the reflection is $u\mapsto5-u$, so the objects $0,1,2$ of the plus
half face $5,4,3$. The base pairs $\{0,3\}$, $\{1,4\}$, $\{2,5\}$ therefore join $a$
with $c'$, $b$ with $b'$ and $c$ with $a'$, and $u\mapsto5-u$ maps $\{0,3\}$ to $\{2,5\}$
and $\{1,4\}$ to itself. At $\iota$ the plus half is the layer $\ell=0$, and the base pair
$\{j,j+3\}$ has one base in each component.
\end{proof}

\subsection{The chamber symmetry for the tilted law}\label{subsec:positive-root-chamber}

This subsection proves the cap of the chamber cut L20 of a positive root. The proof
follows that of \cref{subsec:middle-l20} for the uniform root: the soft kernel is a
swapped rectangle, its leading part is a complete star matrix synthesized by the actual
frame, and everything else is a remainder of order $q$. The tilt adds to the leading
part one exact matrix term, the coupling of the pair features by the linear term
$2\ulaw\Hbar$ (\cref{lem:root-pair-matrix}); its square terms are, up to order $q$,
multiples of the untilted kernel (\cref{lem:root-square}); the
linear term has an explicit remainder (\cref{lem:root-linear}); and the normalization,
the half marginal and the pointwise conditional score mean are bounded in
\cref{lem:root-normalization}. All estimates are written so that they can be evaluated
at a prime endpoint and hold for every larger prime.

\paragraph{The chart.}
At $\vartheta_0$ the half graph of the root is $G_{\rm F,2}$
(\cref{lem:layer12-geometry}(b)), and the half reference is
$\mu_p=\lambda_p(G_{\rm F,2})=\lambda_{F,p}$, the uniform law on injective tuples
$x=(a_0,a_1,b_0,b_1,c_0,c_1)$ with $a_i\to c_j$. We use the notation of
\cref{subsec:middle-l20}: the \emph{core} $C=(a_0,a_1,c_0,c_1)$ and the \emph{pair}
$b=(b_0,b_1)$ of $x$; the independent laws $\mu^{\rm core}$ (the uniform law of the
admissible cores, injective directed copies of $K_{2,2}$) and $\mu^{\rm pair}$ (the uniform
law of the ordered distinct pairs) of \cref{prop:l20-rectangular}; and the features
$V_k$, $W_S$, $\mathbf u,\mathbf v,\mathbf w,\mathbf z,\mathbf e$ of
\cref{lem:l20-star-matrix}, the same formulas defining them on $\mu_p$ and on
$\mu^{\rm core}\otimes\mu^{\rm pair}$. The law $\mu_p$ is $\mu^{\rm core}\otimes\mu^{\rm pair}$
conditioned on $b\cap C=\emptyset$, an event of probability
\[
 h_p=\frac{(1-4q)(1-5q)}{1-q}\ \ge\ h'_p=(1-4q)(1-5q),
\]
and $J_\mu\colon L^2(\mu_p)\to L^2(\mu^{\rm core}\otimes\mu^{\rm pair})$ denotes its
isometric embedding (\cref{lem:restriction}(a)). The rectangle of
\cref{prop:l20-rectangular} is built from $\mathsf K(z,C)$,
$\bar{\mathsf K}(C)=\E_z\mathsf K(z,C)$, $\mathsf R(C;b)=\mathsf K(b_0,C)\mathsf K(b_1,C)$ and
$\mathfrak r=\mathsf R\one$:
\[
 \mathsf R=\mathfrak r\otimes\one+\mathsf R_\star+\mathsf R_\circ,\qquad
 \mathsf R_\star=r\bigl(\sqrt{30}\,\mathbf z\mathbf v^*+2\mathbf w\mathbf e^*\bigr),\qquad
 \|\mathsf R_\circ\|\le N_A,\qquad\mathsf R_\circ\one=\mathsf R_\circ^*\one=0
\]
\eqref{eq:l20-rectangle}, with $\|\mathbf z\|^2\le g_{\mathbf z}$,
$\|\mathbf w\|,\|\mathbf v\|,\|\mathbf e\|\le1$ (\cref{prop:l20-rank4}) and
$\|\mathsf R_\star\|\le k_4r$ (\cref{prop:l20-rectangular}(b)), and $\mathsf R^*\one=\E\mathfrak r$
is constant. The soft L20 kernel
on $L^2(\mu^{\rm core}\otimes\mu^{\rm pair})$ is
$K_{\rm soft}=T[\mathsf R,\mathsf R]=(\mathsf R\otimes\mathsf R^*)\operatorname{Swap}$
\eqref{eq:l20-swap}, with kernel $\mathsf R(C;b')\mathsf R(C';b)$ at a pair of halves
$x=(C,b)$, $x'=(C',b')$. Further
\begin{gather*}
 Z_4=(1-q)(1-3q)(1-7q),\qquad\eta=11r+2q,\qquad c_0=\frac4{1-q}+2q,\\
 \mathfrak r_+=\frac{(1+\eta)^2}{1-q},\qquad\mathfrak r_-=(1-\eta)(1-\eta-16q),
\end{gather*}
so that $\mathfrak r_-\le\mathfrak r\le\mathfrak r_+$, as $|\bar{\mathsf K}-1|\le\eta$
(\cref{prop:l20-rectangular}(d)). We write $\mathfrak r_2=s_p\ge\|\mathfrak r\|_2$ for the
bound of \cref{prop:l20-rectangular}(d), $\mathfrak m_+\ge|\E\bar{\mathsf K}-1|$ for that of
\cref{prop:l20-rank4}, $\Gamma_{20}\ge\|\mathcal F\|^2$ for the frame cap of
\cref{lem:l20-star-matrix}(b) ($\mathcal F$ the synthesis of the eighteen features on
$L^2(\mu_p)$, in the coefficient space of \cref{lem:l20-star-matrix}), $N_A$, $D_{\rm bal}$
and $\eps_Z$ for the caps and the normalizer error of
\cref{prop:l20-rectangular,prop:l20-collisions}, $\chi^+_p\ge\|T_0\one-\one\|^2$ for the
root-half chi-square bound of \cref{subsec:middle-l20} (\cref{lem:record:root-marginal}(ii)),
and $R_p$ for that of \cref{lem:record:root-marginal}(i). For a feature map $Y$ on
$\mu^{\rm core}\otimes\mu^{\rm pair}$ and the same map $Y_\mu$ on $\mu_p$,
$J_\mu^*Y=\sqrt{h_p}\,Y_\mu$ (\cref{lem:restriction}(a)); hence a leading operator
$YMY'^*$ of the soft kernel becomes $Y_\mu MY_\mu'^*$ under the compression
$h_p^{-1}J_\mu^*(\cdot)J_\mu$, with no factor $h_p^{-1}$, which only norm
remainders pay.

The cut operator of the uniform root on $L^2(\mu_p)$ is the operator $T$ of
\cref{subsec:middle-l20}, which we write $T_0$:
\[
 T_0=K_0/Z_0,\qquad K_0=K^\mu_{\rm soft}-T_{\rm coll},\qquad
 K^\mu_{\rm soft}=h_p^{-1}J_\mu^*K_{\rm soft}J_\mu ,
\]
where $K^\mu_{\rm soft}$ is the soft kernel as an operator on $L^2(\mu_p)$, $T_{\rm coll}$ is
its restriction to the pairs of halves with a common label, a nonnegative kernel with
$\|T_{\rm coll}\|\le D_{\rm bal}$ (\cref{prop:l20-collisions}(a)), and
$Z_0=\ip\one{K_0\one}\ge1-\eps_Z$ (\cref{prop:l20-completion,prop:l20-rectangular}); on
cross-injective pairs every factor of the soft kernel is $0$ or $2$, so the \emph{raw
kernel} $K_0$ is $2^{16}$ times the indicator that $(x,x')$ join to a
labelling of $G_{\rm L,2}$. By
\cref{lem:spec:cutkernel}(f) the cut operator of $\Gamma_{p,\ulaw}$ on $L^2(\mu_p)$ is
\[
 T_\ulaw=\frac{T_0\circ w_\ulaw}{\E_{\Gamma_{p,0}}w_\ulaw}
 =\frac{K_0\circ w_\ulaw}{\ip\one{(K_0\circ w_\ulaw)\one}},
\]
self-adjoint because $w_\ulaw$ is invariant under the reflection
(\cref{lem:positive-root-properties}(c)).

\paragraph{The scores as kernels.}
For a pair $y=(y_0,y_1)$ of labels put $\psi_y(z)=\chi(z-y_0)\chi(z-y_1)$. By
\cref{lem:positive-root-properties}(c), at $\vartheta_0$
\[
 \Hs0(x,x')=\frac1r\E_z\psi_a(z)\psi_{c'}(z),\qquad
 \Hs1(x,x')=\frac1r\E_z\psi_b(z)\psi_{b'}(z),\qquad
 \Hs2(x,x')=\frac1r\E_z\psi_c(z)\psi_{a'}(z),
\]
where $a,b,c$ are the pairs of $x$ and $a',b',c'$ those of $x'$; so the multiplier of a
score is a scalar product of rows of raw pair maps, the setting of
\cref{lem:swapped-rectangle}(d), and we call $\Hs0,\Hs2$ the \emph{core scores} and $\Hs1$
the \emph{private score}. Let $\mathcal U_a=W_{01}/r$ and $\mathcal U_c=W_{45}/r$ be the
normalized pair features of the core and $\mathcal U_b=W_{23}/r=\mathbf e$ that of the
private pair. Since $\E_z\psi_y(z)=-q$ for distinct labels, $W_{01}$ has the kernel
$\psi_a(z)+q$ on mean-zero label functions, and
\begin{equation}\label{eq:root-score-feature}
 \Hs0=r\,\mathcal U_a\mathcal U_c^*+qr,\qquad
 \Hs1=r\,\mathcal U_b\mathcal U_b^*+qr,\qquad
 \Hs2=r\,\mathcal U_c\mathcal U_a^*+qr
\end{equation}
as kernels on the half, the term $qr$ denoting the constant kernel
$qr\,\one\otimes\one$. In the coefficient space of \cref{lem:l20-star-matrix} the three couplings
are $e_{01}e_{45}^*$, $e_{23}e_{23}^*$ and $e_{45}e_{01}^*$.

\begin{lemma}[The pair-feature matrix of a tilt]\label{lem:root-pair-matrix}
In the coefficient space of \cref{lem:l20-star-matrix}, with orthonormal coordinates $e_i$
($0\le i\le5$) and $e_S$ (the twelve supports) and $\Phat$ specialized to $i$
(\cref{rem:realified-gates}), let $M$ be the star matrix of the soft L20 kernel and, for
$\lambda\ge0$,
\[
 M_\lambda=M+\tfrac\lambda3\bigl(e_{01}e_{45}^*+e_{45}e_{01}^*+e_{23}e_{23}^*\bigr).
\]
Then $M_\lambda$ is unitarily equivalent to the direct sum of
$\left(\begin{smallmatrix}0&2&\sqrt{26}\\2&\lambda/3&0\\\sqrt{26}&0&0\end{smallmatrix}\right)$,
$\left(\begin{smallmatrix}0&2\\2&\lambda/3\end{smallmatrix}\right)$, $(-\lambda/3)$ and the
zero matrix of size twelve; its characteristic polynomial is
$x^{12}(x+\frac\lambda3)(x^2-\frac\lambda3x-4)(x^3-\frac\lambda3x^2-30x+\frac{26\lambda}3)$.
Consequently, for $0\le\lambda\le3\sqrt{26}$ the least eigenvalue of $M_\lambda$ is
$-g(\lambda)$, where $g(\lambda)\in(\sqrt{26},\sqrt{30}]$ is the root of
\begin{equation}\label{eq:root-cubic}
 3g^3+\lambda g^2-90g-26\lambda=0 ,
\end{equation}
decreasing in $\lambda$, $g(0)=\sqrt{30}$; every other negative eigenvalue is at least
$-\max(2,\lambda/3)$; and $\|M_\lambda\|\le\sqrt{30}+\lambda/3$. The same holds at
$\Phat=-i$.
\end{lemma}

\begin{proof}
(1) As in the proof of \cref{lem:l20-star-matrix}(a) put
$\hat{\mathbf v}=(e_2+e_3)/\sqrt2$, $\hat{\mathbf w}=(e_0+e_1+e_4+e_5)/2$,
$\hat{\mathbf e}=e_{23}$, $\hat{\mathbf u}=(e_0+e_1-e_4-e_5)/2$ and
\[
 \hat{\mathbf z}=\frac1{\sqrt{15}}\Bigl(\sum_S\sigma_Se_S-2i\,\hat{\mathbf u}\Bigr),\qquad
 \sigma_S=(-1)^{|S\cap\{0,1\}|},
\]
the sum running over the eleven supports $S\subseteq\{0,1,4,5\}$ of the core stars, and
$\sigma_S$ being the sign of the star $S$ in the expansion of $\mathsf K$. The vectors
$\hat{\mathbf v},\hat{\mathbf z},\hat{\mathbf w},\hat{\mathbf e}$ are orthonormal
($|\hat{\mathbf z}|^2=(11+4)/15$), and a comparison of entries gives
\[
 M=\sqrt{30}\,(\hat{\mathbf v}\hat{\mathbf z}^*+\hat{\mathbf z}\hat{\mathbf v}^*)
 +2(\hat{\mathbf w}\hat{\mathbf e}^*+\hat{\mathbf e}\hat{\mathbf w}^*).
\]

(2) $\sigma_{01}=\sigma_{45}=1$, so the coordinates of $\hat{\mathbf z}$ along $e_{01}$ and $e_{45}$
are both $1/\sqrt{15}$. With $\hat{\mathbf u}_\pm=(e_{01}\pm e_{45})/\sqrt2$ this gives
\[
 \hat{\mathbf z}=\sqrt{2/15}\;\hat{\mathbf u}_++\sqrt{13/15}\;\hat{\mathbf z}_\perp
\]
with a unit vector $\hat{\mathbf z}_\perp$ orthogonal to $\hat{\mathbf v}$, $\hat{\mathbf w}$,
$\hat{\mathbf e}$, $\hat{\mathbf u}_+$ and $\hat{\mathbf u}_-$; the six vectors
$\hat{\mathbf v},\hat{\mathbf u}_+,\hat{\mathbf z}_\perp,\hat{\mathbf w},\hat{\mathbf e},
\hat{\mathbf u}_-$ are orthonormal.

(3) $e_{01}e_{45}^*+e_{45}e_{01}^*=\hat{\mathbf u}_+\hat{\mathbf u}_+^*-\hat{\mathbf u}_-\hat{\mathbf u}_-^*$, so the
pair term is $\frac\lambda3(\hat{\mathbf u}_+\hat{\mathbf u}_+^*-\hat{\mathbf u}_-\hat{\mathbf u}_-^*+\hat{\mathbf e}\hat{\mathbf e}^*)$.

(4) As $\sqrt{30}\sqrt{2/15}=2$ and $\sqrt{30}\sqrt{13/15}=\sqrt{26}$,
\begin{multline*}
 M_\lambda=2(\hat{\mathbf v}\hat{\mathbf u}_+^*+\hat{\mathbf u}_+\hat{\mathbf v}^*)
 +\sqrt{26}\,(\hat{\mathbf v}\hat{\mathbf z}_\perp^*+\hat{\mathbf z}_\perp\hat{\mathbf v}^*)
 +\tfrac\lambda3\hat{\mathbf u}_+\hat{\mathbf u}_+^*\\
 +2(\hat{\mathbf w}\hat{\mathbf e}^*+\hat{\mathbf e}\hat{\mathbf w}^*)+\tfrac\lambda3\hat{\mathbf e}\hat{\mathbf e}^*
 -\tfrac\lambda3\hat{\mathbf u}_-\hat{\mathbf u}_-^*.
\end{multline*}
In the orthonormal bases $(\hat{\mathbf v},\hat{\mathbf u}_+,\hat{\mathbf z}_\perp)$, $(\hat{\mathbf w},\hat{\mathbf e})$,
$\hat{\mathbf u}_-$, completed by an orthonormal basis of the complement, on which $M_\lambda$
vanishes, this is the displayed direct sum. Its blocks have the characteristic
polynomials $x^3-\frac\lambda3x^2-30x+\frac{26\lambda}3$, $x^2-\frac\lambda3x-4$ and
$x+\frac\lambda3$.

(5) For $g>0$ the value of the cubic at $x=-g$ is $-f(g)/3$ with
$f(g)=3g^3+\lambda g^2-90g-26\lambda$. Since $f(0)=-26\lambda\le0$ and $f''\ge0$ on
$[0,\infty)$, $f$ has exactly one positive root $g(\lambda)$ (for $\lambda=0$ it is
$\sqrt{30}$); as $f(\sqrt{26})=-12\sqrt{26}<0$ and $f(\sqrt{30})=4\lambda\ge0$, it lies in
$(\sqrt{26},\sqrt{30}]$. At the root $\partial_gf=9g^2+2\lambda g-90>0$ and
$\partial_\lambda f=g^2-26>0$, so $g$ decreases in $\lambda$. So the $3\times3$ block has
the single negative eigenvalue $-g(\lambda)$. The $2\times2$ block has the negative
eigenvalue $\lambda/6-(\lambda^2/36+4)^{1/2}\in[-2,0)$, and the last block $-\lambda/3$;
for $\lambda\le3\sqrt{26}<3g(\lambda)$ both lie above $-g(\lambda)$. The $3\times3$ block is
$B_0+\frac\lambda3\hat{\mathbf u}_+\hat{\mathbf u}_+^*$, where $B_0$ has the characteristic polynomial
$x^3-30x$, so its norm is at most $\sqrt{30}+\lambda/3$; the $2\times2$ block has norm
$\lambda/6+(\lambda^2/36+4)^{1/2}\le2+\lambda/3$. At $\Phat=-i$ the matrix is the
complex conjugate (\cref{rem:realified-gates}), with the same spectrum.
\end{proof}

The orientation of the pair term matters. The scores couple the core pair $a$ of one
half with the core pair $c'$ of the other, and the leading vector $\hat{\mathbf z}$ has a
component along $\hat{\mathbf u}_+$, on which this coupling is positive. With the opposite sign
on the two core couplings, $\hat{\mathbf u}_+$ and $\hat{\mathbf u}_-$ exchange their roles, the
cubic becomes $3g^3-\lambda g^2-90g+26\lambda=0$, and its root exceeds $\sqrt{30}$.

\begin{lemma}[Square of a score on the rectangle]\label{lem:root-square}
Suppose $k_4r+N_A<1$. Then for each score, on $L^2(\mu^{\rm core}\otimes\mu^{\rm pair})$,
\[
 \bigl\|K_{\rm soft}\circ\Hs j^2-(1-2q)^2\,K_{\rm soft}\bigr\|\le256c_0\mathfrak r_+q,
 \qquad
 \bigl\|K_{\rm soft}\circ\Hs j^2-K_{\rm soft}\bigr\|\le4q\mathfrak r_2^2+256c_0\mathfrak r_+q .
\]
\end{lemma}

\begin{proof}
Take the core score $\Hs0$; $\Hs2$ is its transpose. For a pair $y$ of distinct labels let
$\Psi_y(z,z')=\psi_y(z)\psi_y(z')$ in $L^2(\Fp^2)$, so that
$\Hs0^2(x,x')=q^{-1}\ip{\Psi_a}{\Psi_{c'}}$. By \cref{cor:raw-pair-maps}(b),
$\Psi_y=r(1-2q)d+\rho_y$, where $d=\sqrt p\,\one_{z=z'}$ is a unit vector, $\rho_y\perp d$
is the row of the remainder of \eqref{eq:raw-pair-split}, and the map
$F\mapsto\ip{\rho_y}F$ has norm at most $\sqrt{c_0}\,q$ into $L^2$ of a uniform distinct
pair $y$. Under $\mu^{\rm core}$ the pairs $a$ and $c$ of a core are uniform distinct
pairs, so the same bound holds into $L^2(\mu^{\rm core})$. Since $\rho\perp d$, the cross terms vanish and
\[
 \Hs0^2=(1-2q)^2+q^{-1}\ip{\rho_a}{\rho_{c'}},\qquad
 K_{\rm soft}\circ\Hs0^2-(1-2q)^2K_{\rm soft}
 =q^{-1}\,T[\mathsf R,\mathsf R]\circ\ip{\rho_a}{\rho_{c'}} .
\]
By \cref{lem:swapped-rectangle}(d) the right side is
$q^{-1}(\mathcal Y_a\otimes I)\operatorname{Swap}(\mathcal Y_c^*\otimes I)$, the auxiliary
label being shared and the private variables exchanged, where $\mathcal Y_a$ has the kernel
$\mathsf R(C;b)\rho_{a(C)}$ from $L^2(\mu^{\rm pair})\otimes L^2(\Fp^2)$ to
$L^2(\mu^{\rm core})$, and $\mathcal Y_c$
likewise. The Gram $\mathcal Y_a\mathcal Y_a^*$ is the pointwise product of the positive
semidefinite kernel $\mathsf R\mathsf R^*$, whose diagonal is
$\E_b\mathsf R(C;b)^2=256\,\mathfrak r(C)\le256\,\mathfrak r_+$ ($\mathsf R$ takes the values
$0$ and $256$), with the Gram of the residual pair map, which lies between $0$ and
$c_0q^2$. By the positive Schur order of \cref{lem:swapped-rectangle}(b),
$\|\mathcal Y_a\|^2\le256c_0\mathfrak r_+q^2$, and likewise for $\mathcal Y_c$; this proves the
first bound for the core scores. For the private score the pairs $b$ and $b'$ are the input
variables of $\mathsf R$: in the variables $(b,C)$ the soft kernel is
$T[\mathsf R^*,\mathsf R^*]$ (the last sentence of \cref{lem:swapped-rectangle}), and the same
argument applies with the rows of $\mathsf R^*$, whose square norm is the constant
$256\,\E\mathfrak r\le256\,\mathfrak r_+$.

For the second bound, $\mathsf R=\mathfrak r\otimes\one+\mathsf R_c$ with $\mathsf R_c\one=0$
and $\mathsf R_c^*\mathfrak r=0$ (\cref{prop:l20-rectangular}(a)), so $\mathsf R$ maps $\one$
to $\mathfrak r$ and $\one^\perp$ into $\mathfrak r^\perp$, and its norm is the larger of
$\|\mathfrak r\|_2\le\mathfrak r_2$ and $\|\mathsf R_c\|\le k_4r+N_A<1$
(\cref{prop:l20-rectangular}(b),(d)). As $\mathfrak r_2=s_p\ge1/(1-q)>1$,
$\|K_{\rm soft}\|=\|\mathsf R\|^2\le\mathfrak r_2^2$ (\cref{lem:swapped-rectangle}(a)), and
$1-(1-2q)^2\le4q$.
\end{proof}

The estimate bounds the square of a score on the whole rectangle at once,
with the coefficient $4\mathfrak r_2^2+256c_0\mathfrak r_+$ of $q$, which tends to
$1028$ as $p\to\infty$. Bounded class by class instead, the remainder class
of the rectangle alone would give about $N_A^2/q$, about $8\cdot10^4q$.

\paragraph{Linear prices.}
The linear part of the tilt is $K_{\rm soft}\circ\Hbar$. Its finite price uses the
following constants. Put
\[
 d_{\mathsf K}=\Bigl(\frac{5q+34qr+11q^3}{Z_4}\Bigr)^{1/2}+2q,\quad
 d_{\mathfrak r}=\frac{(2+\eta+16q)d_{\mathsf K}+15q}{1-q},\quad
 \delta_{\mathfrak r}=\frac{d_{\mathfrak r}}{\sqrt{1-q}},\quad
 A_{\mathfrak r}=1+\delta_{\mathfrak r}.
\]
In the expansion of $\bar{\mathsf K}-1$ the single terms vanish and the pair terms sum to
exactly $2q$; the higher terms are $\mathcal W/p$ in the notation of
\cref{lem:record:weighted-moment} (signs $(-1,-1,1,1)$ on $(a_0,a_1,c_0,c_1)$), and the core
law $\mu^{\rm core}$ has the density $P_{\rm adm}D_4/Z_4$ relative to independent uniform
labels ($D_4$ is the indicator that the four labels are distinct, and $P_{\rm adm}$ is $16$
on the admissible cores, of total probability $Z_4/16$, and $0$ on the other distinct
quadruples), so
$\E_{\mu^{\rm core}}(\mathcal W/p)^2\le\E[P_{\rm adm}\mathcal W^2]/(p^2Z_4)$. Hence
\cref{lem:record:weighted-moment}(b) gives $\|\bar{\mathsf K}-1\|_2\le d_{\mathsf K}$, and since
$(1-q)(\mathfrak r-1)=(\bar{\mathsf K}-1)(\bar{\mathsf K}+1-16q)-15q$,
$\|\mathfrak r-1\|_2\le d_{\mathfrak r}$. For a core function $f$ invariant under the
square-affine maps and a normalized core star $\mathcal U_S=W_S/r$, the Gram of
$f\mathcal U_S$ commutes with the affine group and is self-adjoint, hence scalar
(\cref{lem:spec:affpair}(d)), and its trace gives $\|f\mathcal U_S\|\le\|f\|_2/\sqrt{1-q}$
(as in \cref{lem:general-star}(a)). So multiplying a core star by $\mathfrak r$ changes it
by at most $\delta_{\mathfrak r}$ and enlarges its norm by at most the factor
$A_{\mathfrak r}$; for the private features, which are independent of the core, the change
is at most $d_{\mathfrak r}\le\delta_{\mathfrak r}$.

\emph{Product maps.} Fix a same-side core pair $A_0$ (the pair $a$ or the pair $c$), with
normalized pair feature $\mathcal U_{A_0}$. For a core support $S$ with $2\le|S|\le4$ the
product map $\mathcal U_S\times\mathcal U_{A_0}\colon\mathcal H\otimes\mathcal H\to L^2(\mu^{\rm core})$,
$F\mapsto\E_{t,z}\,\mathcal U_S(\cdot;t)\,\mathcal U_{A_0}(\cdot;z)F(t,z)$
($\mathcal H=L^2_0(\Fp)$), is bounded on three orthogonal pieces of $\mathcal H\otimes\mathcal H$: the functions
vanishing on the diagonal $t=z$, the mean-zero diagonal functions
$D_{\mathcal H}f=(P_0\otimes P_0)\sqrt p\,\one_{t=z}f(t)$, $f\in\mathcal H$, whose Gram is $(1-2q)I$, and the
constant diagonal vector $e_0=(p\one_{t=z}-1)/\sqrt{p-1}$ of
\cref{lem:swapped-rectangle}(c). With $d=|S\triangle A_0|$, $c=|S\cap A_0|$,
$u=|S\cup A_0|$, the densities $\mathfrak d_2=1/(1-q)$, $\mathfrak d_3=4/((1-q)(1-3q))$,
$\mathfrak d_4=16/Z_4$ of a same-side pair, a triple of core labels
consisting of one same-side pair and one further core label, and the whole
core relative to independent uniform labels, and the bounds $\omega_1=0$, $\omega_2=q$,
$\omega_3=2r$, $\omega_4=2r+q$ for the complete character mean $|\E_t\psi_S(t)|$ of
$1,\dots,4$ distinct labels, put
\begin{gather*}
 \varrho^o_S=\begin{cases}2\sqrt{\mathfrak d_2},&S=A_0,\\\sqrt{\mathfrak d_u},&S\ne A_0,\end{cases}
 \qquad
 \varrho^d_S=\frac{\one_{d\ne0}+cr+q+\omega_{|S|}}{\sqrt{1-2q}},\qquad
 \varrho^c_S=\begin{cases}0,&S=A_0,\\[2pt]
 \dfrac{\omega_d+cq+\omega_{|S|}q}{r\sqrt{1-q}},&S\ne A_0,\end{cases}\\
 R_S=\bigl((\varrho^o_S)^2+(\varrho^d_S)^2+(\varrho^c_S)^2\bigr)^{1/2}.
\end{gather*}
Then $\mathcal U_S\times\mathcal U_{A_0}$ has norm at most $R_S$, except that for $S=A_0$
the vector $e_0$ has the constant image $h_{A_0}\one$,
$h_{A_0}=(1-2q-q^2)/(r\sqrt{1-q})$, which is set apart. Indeed, on the off-diagonal piece
choose an output label in $S\setminus A_0$ and one in $A_0$: with the two auxiliary labels
they span a matching or a three-edge path, of norm $q$ by two Paley transforms
(\cref{lem:matching:kernel}(a)), the other labels entering as bounded multipliers and the
density of the $u$ labels costing $\sqrt{\mathfrak d_u}$; for $S=A_0$ the incidence is a
$K_{2,2}$, of counting norm $2p$ (\cref{thm:dl:K}), whence $2q$; the separate distinctness
compressions stay in place. On $D_{\mathcal H}\mathcal H$ the image is
$r^{-1}(W_{S\triangle A_0}-\text{$c$ holes}-w_{A_0}W_S-w_SW_{A_0})$, with $w_{A_0}=-q$, each
hole (evaluation at a common label) costing $q$ and every nonempty star having norm at
most $r$. On $e_0$ the image is $(\E_t\psi_S\psi_{A_0}-w_Sw_{A_0})/(r\sqrt{1-q})$, bounded
by $\varrho^c_S$ for $S\ne A_0$ and equal to $h_{A_0}$ for $S=A_0$. A singleton inside
$A_0$ times $\mathcal U_{A_0}$ costs $1/\sqrt{1-q}$ and one outside $A_0$ costs
$\sqrt{\mathfrak d_3}$; so the products of $\mathbf w$, and of $\mathbf u$ and
$\mathbf u\Phat$ ($\Phat$ is unitary on $\mathcal H$), with $\mathcal U_{A_0}$ have norm at most
$R_{\mathbf w}=1/\sqrt{1-q}+\sqrt{\mathfrak d_3}$. With
$\mathbf z=(\sum_S\sigma_S\mathcal U_S-2\mathbf u\Phat)/\sqrt{15}$ and $\sigma_{A_0}=1$,
\begin{gather*}
 \mathbf z\times\mathcal U_{A_0}=h'_{\mathbf z}\,\one e_0^*+(\text{residual}),\qquad
 0\le h'_{\mathbf z}\le h_{\mathbf z}=\frac1{r\sqrt{15(1-q)}},\\
 \|\text{residual}\|\le R_{\mathbf z}=\frac{\sum_SR_S+2R_{\mathbf w}}{\sqrt{15}},
\end{gather*}
the sum running over the eleven core supports. For the private pair, $\mathbf v\times\mathbf e$
has norm at most $R_{\mathbf v}=(2/(1-q))^{1/2}$, and
$\mathbf e\times\mathbf e=h'_{\mathbf e}\one e_0^*+(\text{residual})$ with residual norm at most
$R_{\mathbf e}=(4/(1-q)+(2r+2q)^2/(1-2q))^{1/2}$ and
$0\le h'_{\mathbf e}\le h_{\mathbf e}=1/(r\sqrt{1-q})$, by the computation for $S=A_0$, the
private pair being a uniform distinct pair. As $p\to\infty$, $R_{\mathbf z}$ tends to about
$10.14$, $R_{\mathbf w}$ to $3$, $R_{\mathbf v}$ to $\sqrt2$ and $R_{\mathbf e}$ to $2$. The class
prices are
\begin{align*}
 \eps_{00}&=r(2\delta_{\mathfrak r}+\delta_{\mathfrak r}^2),\qquad
 \eps_{01,c}=2qA_{\mathfrak r}(\sqrt{30}R_{\mathbf z}+2R_{\mathbf w}),\\
 \eps_{01,b}&=2q\bigl\{\sqrt{30}\sqrt{g_{\mathbf z}}\,\mathfrak r_2R_{\mathbf v}
   +2(\mathfrak r_2R_{\mathbf e}+h_{\mathbf e}d_{\mathfrak r})\bigr\},\\
 \eps_{11,c}&=r^3\bigl\{30\bigl(R_{\mathbf z}^2+2h_{\mathbf z}R_{\mathbf z}
   +\tfrac{h_{\mathbf z}^2}{p-1}\bigr)+4\sqrt{30}(R_{\mathbf z}+h_{\mathbf z})R_{\mathbf w}
   +4R_{\mathbf w}^2\bigr\},\\
 \eps_{11,b}&=r^3\bigl\{30g_{\mathbf z}R_{\mathbf v}^2
   +4\sqrt{30}\sqrt{g_{\mathbf z}}\,R_{\mathbf v}(R_{\mathbf e}+h_{\mathbf e})
   +4\bigl(R_{\mathbf e}^2+2h_{\mathbf e}R_{\mathbf e}+\tfrac{h_{\mathbf e}^2}{p-1}\bigr)\bigr\},\\
 \eps_{02}&=2A_{\mathfrak r}N_A,\qquad
 \eps_{12,c}=2N_A\bigl\{\sqrt{30}(rh_{\mathbf z}+rR_{\mathbf z})+2rR_{\mathbf w}\bigr\},\\
 \eps_{12,b}&=2N_A\bigl\{\sqrt{30}\sqrt{g_{\mathbf z}}\,rR_{\mathbf v}
   +2(rh_{\mathbf e}+rR_{\mathbf e})\bigr\},\qquad
 \eps_{22}=N_A^2/r,
\end{align*}
and, a core value counted twice and the private value once,
\begin{gather*}
 \eps_{\rm lin}=\eps_{00}+\frac{2\eps_{01,c}+\eps_{01,b}}3+\frac{2\eps_{11,c}+\eps_{11,b}}3
 +\eps_{02}+\frac{2\eps_{12,c}+\eps_{12,b}}3+\eps_{22}+qr\mathfrak r_2^2,\\
 \mathfrak l_0=\frac{2\sqrt{30}\,rh_{\mathbf z}A_{\mathfrak r}+2rh_{\mathbf e}\mathfrak r_2}3 .
\end{gather*}

\begin{lemma}[Linear root trace]\label{lem:root-linear}
Suppose $k_4r+N_A<1$. On $L^2(\mu^{\rm core}\otimes\mu^{\rm pair})$,
\[
 P_0(K_{\rm soft}\circ\Hbar)P_0=\frac r3\bigl(\mathcal U_a\mathcal U_c^*+\mathcal U_c\mathcal U_a^*
 +\mathcal U_b\mathcal U_b^*\bigr)+\mathcal R_{\rm lin},\qquad
 \|\mathcal R_{\rm lin}\|\le\eps_{\rm lin},
\]
the leading term being the pair term of \cref{lem:root-pair-matrix} with coefficient
$r/3$ per score, synthesized by the features. Moreover
$\|P_0(K_{\rm soft}\circ\Hbar)\one\|\le\mathfrak l_0r+\eps_{\rm lin}$ and
$\|P_0(K_{\rm soft}\circ\Hbar)\|\le r(1+2\mathfrak l_0)+\eps_{\rm lin}$, and
$\mathfrak l_0\to(2\sqrt2+2)/3$.
\end{lemma}

\begin{proof}
By \eqref{eq:root-score-feature}, $K_{\rm soft}\circ\Hs0=r\,K_{\rm soft}\circ\mathcal U_a\mathcal U_c^*
+qr\,K_{\rm soft}$, and the last term has norm at most $qr\mathfrak r_2^2$ (proof of
\cref{lem:root-square}); the same holds for the other scores. Expand both factors of
$K_{\rm soft}=T[\mathsf R,\mathsf R]$ by \eqref{eq:l20-rectangle}: $K_{\rm soft}$ is the
sum of the operators $T[X,Y]$ with $X,Y$ the rank $\mathfrak r\otimes\one$, the star
$\mathsf R_\star$ or the remainder $\mathsf R_\circ$, six classes up to transposition. For a
core score write the kernel of $T[X,Y]\circ\mathcal U_a\mathcal U_c^*$ at $(x,x')$ as
$\E_z[X(C;b')\mathcal U_a(C;z)]\,[Y(C';b)\mathcal U_c(C';z)]$: it factors through the
shared auxiliary label $z$ and the exchange of the private variables
(\cref{lem:swapped-rectangle}(d)), and its norm is at most the product of the norms of
the two factors, each a product map of the arms (core or pair feature
factors) of $X$ or $Y$ with a pair feature. The
star $\mathsf R_\star$ has the two arms $(\mathbf z,\mathbf v)$ and $(\mathbf w,\mathbf e)$
with the coefficients $\sqrt{30}\,r$ and $2r$.

\emph{Rank/rank.} The factors are $\mathfrak r\,\mathcal U_a$ and $\mathfrak r\,\mathcal U_c$,
so the class is $r\,(\mathfrak r\,\mathcal U_a)(\mathfrak r\,\mathcal U_c)^*$. Since
$P_0(\mathfrak r\,\mathcal U_a)-\mathcal U_a=P_0((\mathfrak r-1)\mathcal U_a)$ has norm at most
$\delta_{\mathfrak r}$, it is the leading term $r\,\mathcal U_a\mathcal U_c^*$ up to
$\eps_{00}$ after centring on both sides.

\emph{Rank/star.} The factors are $\mathfrak r\,\mathcal U_a$ (norm at most $A_{\mathfrak r}$)
tensored with the private arm, and the product of the core arm with $\mathcal U_c$, with the
coefficient $r\cdot r=q$ of the star and the score. By the product maps the second factor is
bounded by $R_{\mathbf z}$ or $R_{\mathbf w}$ except for the large constant
$h'_{\mathbf z}\one e_0^*$ of $\mathbf z\times\mathcal U_c$, which yields an operator whose
input is the constant function of $x'$: a constant-to-feature rank, killed by $P_0$ on the
right. The two placements give $\eps_{01,c}$. For the private score the large constant is
that of $\mathbf e\times\mathbf e$; its output is $\mathfrak r$ times a constant, and
$P_0(\mathfrak r h)=hP_0\mathfrak r$ is paid by $h_{\mathbf e}d_{\mathfrak r}$ in $\eps_{01,b}$.

\emph{Star/star.} The coefficient is $r\cdot r\cdot r$ and the four arm pairs give the
products of the bounds above. The only term with a large constant on both sides is
$\mathbf z\times\mathcal U_a$ against $\mathbf z\times\mathcal U_c$ (for the private score
$\mathbf e\times\mathbf e$ against itself); there the two vectors $e_0$ meet through the
exchange of the auxiliary labels, and the entangled swap of
\cref{lem:swapped-rectangle}(c) reduces the term to $I/(p-1)$. This gives $\eps_{11,c}$ and
$\eps_{11,b}$.

\emph{Classes with the remainder.} Apply $\mathsf R_\circ$ fibrewise in its own core and
private variables, keeping the shared auxiliary label (it is not duplicated), then the free
feature on the other side: rank/remainder costs $2A_{\mathfrak r}N_A$, star/remainder
$\eps_{12,c}$ or $\eps_{12,b}$ (here the large constant costs $rh_{\mathbf z}$ or
$rh_{\mathbf e}$), and remainder/remainder $N_A^2/r$, the raw pair rows having norm at most
one.

Averaging the three scores, a core value twice and the private value once, gives the
first assertion. The constant-to-feature ranks are explicit: at a core score the vector
obtained by contracting $e_0$ with $\mathfrak r\,\mathcal U\otimes\mathbf v$ has norm at most
$A_{\mathfrak r}$ and mean zero, by translation invariance and independence, with the
coefficient $\sqrt{30}\,rh_{\mathbf z}\cdot r$; at the private score the vector from
$\mathbf w\otimes\mathbf e$ has norm at most $\E\mathfrak r\le\mathfrak r_2$, with the
coefficient $2rh_{\mathbf e}\cdot r$. Exactly one of the two orientations survives in
$P_0(K_{\rm soft}\circ\Hbar)\one$, which gives $\mathfrak l_0r$ rather than twice that
number; the other parts of $P_0(K_{\rm soft}\circ\Hbar)\one$ are bounded by their norms. The
leading pair term has norm at most $r$, which gives the bound for
$\|P_0(K_{\rm soft}\circ\Hbar)\|$. As $p\to\infty$, $A_{\mathfrak r},\mathfrak r_2\to1$ and
$\mathfrak l_0\to(2\sqrt{30}/\sqrt{15}+2)/3$.
\end{proof}

\begin{lemma}[Normalization, marginal and conditional prices]\label{lem:root-normalization}
Let $p\ge2^{20}$. Put
\begin{gather*}
 \bar{\mathfrak r}=\frac{1+2\mathfrak m_++d_{\mathsf K}^2}{1-q},\qquad
 \mathfrak s_b=\frac{512\,\bar{\mathfrak r}q^2}{1-q},\qquad
 \mathfrak s_c=\frac{32\,\mathfrak r_+(1+\eta)^2q^2}{Z_4},\qquad
 \mathfrak s=\frac{2\mathfrak s_c+\mathfrak s_b}3,\\
 c_{\rm lin}=(2+r)D_{\rm bal}+\frac{\mathfrak s}{h'_pr},\qquad
 \eps^{\rm raw}_{\rm lin}=\frac{\eps_{\rm lin}}{h'_p}+c_{\rm lin},\\
 \eps^{\rm raw}_{\rm sq}=\frac{4q\mathfrak r_2^2+256c_0\mathfrak r_+q}{h'_p}
 +\bigl((2+r)^2-1\bigr)D_{\rm bal}+\frac{\mathfrak s}{h'_pq},\qquad
 \eps_{\rm sq}=\frac{\eps^{\rm raw}_{\rm sq}}{1-\eps_Z},\\
 \mathfrak l=\frac1{1-\eps_Z}\Bigl\{\frac{\mathfrak l_0r+\eps_{\rm lin}}{\sqrt{h'_p}}
 +\frac{[r(1+2\mathfrak l_0)+\eps_{\rm lin}]\sqrt{8q/(1-q)}}{h'_p}+c_{\rm lin}\Bigr\},\\
 h_0=\frac{72q+37qr+2qr/(1-q)}{1-36q}.
\end{gather*}
\begin{enumerate}[label=(\alph*)]
\item If $k_4r+N_A<1$, then in the fixed half chart $L^2(\mu_p)$,
 $\bigl\|K_0\circ\tfrac13\sum_j\Hs j^2-K_0\bigr\|\le\eps^{\rm raw}_{\rm sq}$ and
 \[
  \Bigl\|P_0(K_0\circ\Hbar)P_0-\frac r3\,\mathcal F\bigl(e_{01}e_{45}^*+e_{45}e_{01}^*
  +e_{23}e_{23}^*\bigr)\mathcal F^*\Bigr\|\le\eps^{\rm raw}_{\rm lin};
 \]
 consequently $\bigl\|T_0\circ\tfrac13\sum_j\Hs j^2-T_0\bigr\|\le\eps_{\rm sq}$ and
 $\|P_0(T_0\circ\Hbar)\one\|\le\mathfrak l$.
\item For each score, $|\E_{\Gamma_{p,0}}\Hs j|\le h_0$ and
 $\E_{\Gamma_{p,0}}\Hs j^2\ge1-4q-36q(2+r)^2-2q/(1-q)$; in particular
 $|\E_{\Gamma_{p,0}}\Hbar|\le h_0$.
\item Put $\Delta_{\rm core}=(1-q)^4-16r-144q-512qr-512q^2$, positive for $p\ge2^{20}$,
 and
 \begin{gather*}
  c_{\rm core}=\frac{148q+616qr+580q^2}{r\Delta_{\rm core}}+2r,\qquad
  \varphi_b=\frac{24r+8q}{1-\eta},\qquad d_b=\frac{16q}{1-\eta-16q},\\
  c_b=\frac{4q+(1+d_b)\varphi_b^2+d_b}r,\qquad c_{\rm avg}=\frac{2c_{\rm core}+c_b}3,\qquad
  I_{\rm in}=\frac{128\,q\,\mathfrak r_+\bar{\mathfrak r}}{1-\eta},\qquad
  \rho_{\rm in}=\mathfrak r_-^2-I_{\rm in}-D_{\rm bal}.
 \end{gather*}
 If $\rho_{\rm in}>0$, then for every half $x$
 \[
  \bigl|\E_{\Gamma_{p,0}}[\Hbar\mid x]\bigr|\le c_{\rm cond}
  :=\frac{\mathfrak r_+\bar{\mathfrak r}\,c_{\rm avg}+(2+r)(I_{\rm in}+D_{\rm bal})
  +\mathfrak s/r}{\rho_{\rm in}} .
 \]
\end{enumerate}
\end{lemma}

\begin{proof}
(a) \emph{Spikes.} $\bar{\mathfrak r}$ bounds
$\E\mathfrak r\le\E\bar{\mathsf K}^2/(1-q)$, as $\E\bar{\mathsf K}^2\le1+2\mathfrak m_++d_{\mathsf K}^2$. When the two halves share a label, the only large values of a score
occur when the two selected pairs are equal as sets; there
$\Hs j=(1-2q)/r$, and otherwise $|\Hs j|\le2+r$. For the private score, $\{b'\}=\{b\}$ has
probability $2q^2/(1-q)$ and $\mathsf R\le256$; integrating the core of the other half gives
the column mass $\E\mathfrak r\le\bar{\mathfrak r}$, so the rows and columns of the soft kernel
on this event are at most $\mathfrak s_b$. For a core score, $\{c'\}=\{a\}$ has the same
probability; given it, the labels $a'$ are a uniform ordered distinct pair among the
$(p-3)/4$ common predecessors of the pinned pair (\cref{lem:paley-triple}(a)), and both must
lie in the set of labels that extend $(b,c')$ cyclically, of at most $p(1+\eta)/16$ points,
so the conditional mean of $\mathsf R(C';b)$ is at most $16(1+\eta)^2/((1-3q)(1-7q))$;
integrating the pair of the other half gives $\mathfrak r(C)\le\mathfrak r_+$, and this
gives $\mathfrak s_c$. Away from these events a score has
$|\Hs j|\le2+r$, as a single repeated label leaves a character sum of degree two, so
$|\Hs j^2-1|\le(2+r)^2-1$, and the nonnegative kernel $T_{\rm coll}$ contributes in the
displayed directions. The spike prices are row and column bounds of nonnegative kernels
(Schur's test), compressed with the factor $1/h_p\le1/h'_p$.

\emph{Squares.} For $V=\frac13\sum_j\Hs j^2$,
$K_0\circ V-K_0=h_p^{-1}J_\mu^*(K_{\rm soft}\circ V-K_{\rm soft})J_\mu
-T_{\rm coll}\circ(V-1)$, and \cref{lem:root-square,lem:restriction}(b) bound the first term
by $(4q\mathfrak r_2^2+256c_0\mathfrak r_+q)/h'_p$; with the collision bounds this is
$\eps^{\rm raw}_{\rm sq}$, and division by $Z_0\ge1-\eps_Z$ gives $\eps_{\rm sq}$.

\emph{Linear term.} By \cref{lem:restriction}(a) the embedding maps centred functions to
centred functions and $J_\mu^*\one=\sqrt{h_p}\one$, so
$P_0(K^\mu_{\rm soft}\circ\Hbar)P_0=h_p^{-1}P_0J_\mu^*\bigl[P_0(K_{\rm soft}\circ\Hbar)P_0\bigr]
J_\mu P_0$. By \cref{lem:root-linear} and the identity $J_\mu^*Y=\sqrt{h_p}Y_\mu$ for
the features, the leading term becomes
$\frac r3\mathcal F(e_{01}e_{45}^*+e_{45}e_{01}^*+e_{23}e_{23}^*)\mathcal F^*$ and the remainder
has norm at most $\eps_{\rm lin}/h'_p$; the collision part $T_{\rm coll}\circ\Hbar$ has norm at
most $c_{\rm lin}$. For the marginal, \cref{lem:restriction}(b) with both events of
probability $h_p$ gives
\[
 \|P_0(K^\mu_{\rm soft}\circ\Hbar)\one\|\le\frac{\|P_0(K_{\rm soft}\circ\Hbar)\one\|}{\sqrt{h_p}}
 +\|P_0(K_{\rm soft}\circ\Hbar)P_0\|\,\frac{(1-h_p)^{1/2}}{h_p},
\]
where $1-h_p=(8q-20q^2)/(1-q)\le8q/(1-q)$ and \cref{lem:root-linear} bounds the two norms.
Adding $c_{\rm lin}$ and dividing by $Z_0$ gives $\mathfrak l$.

(b) Before global injectivity the root is the product of the uniform laws of its two
components, each a cyclic $K_{2,2,2}$ on the bases of one parity (the
parts being the two layer labels at each of the three bases, with all
edges directed cyclically to the next base), and each score pairs the
two labels of one base of one component with those of a base of the other. Each such pair
is a uniform ordered distinct pair: square-affine maps act transitively on each of the
two orbits of ordered distinct pairs and the layer swap exchanges them. Hence, with
$\E\psi_y(z)=-q$, $\E\Hs j=r^{-1}\E_z(\E\psi_y(z))^2=qr$, and
\[
 \E\Hs j^2=q^{-1}\E_{z,z'}\bigl(\E_y\psi_y(z)\psi_y(z')\bigr)^2
 =(1-2q)^2+\frac{q(1-3q)^2}{1-q}\ \ge\ 1-4q,
\]
as $\E_y\psi_y(z)^2=1-2q$ and $\E_y\psi_y(z)\psi_y(z')=-q(1-3q)/(1-q)$ for $z\ne z'$. The
thirty-six coincidences between the labels of the two components have total probability
at most $36q$; on them $|\Hs j|\le2+r$, except on the events where the two selected pairs
are equal as sets, of probability at most $2q^2/(1-q)$, where $|\Hs j|\le1/r$. Conditioning
on injectivity therefore gives
$|\E_{\Gamma_{p,0}}\Hs j|\le(qr+36q(2+r)+2qr/(1-q))/(1-36q)=h_0$ and
$\E_{\Gamma_{p,0}}\Hs j^2\ge\E\Hs j^2-36q(2+r)^2-2q(1-2q)^2/(1-q)$, whence the bounds.

(c) Fix $x=(C,b)$. On $\mu^{\rm core}\otimes\mu^{\rm pair}$ the soft conditional law of $x'=(C',b')$ given $x$ is
proportional to $\mathsf R(C;b')\mathsf R(C';b)$; it is the product of a law of $C'$ that
depends on $b$ and a law of $b'$ that depends on $C$, and
$\E_{x'}K_{\rm soft}(x,x')=\mathfrak r(C)\E\mathfrak r$.

\emph{A core score.} Take $\Hs0$, which tests the pair $c'$ of $C'$:
$\E[\Hs0\mid x]=r^{-1}\E_z\psi_a(z)\,\E[\psi_{c'}(z)\mid b]$. In independent uniform labels
$(a'_0,a'_1,c'_0,c'_1)$ the law of $C'$ given $b$ has a density proportional to the
product of the four internal edge factors $1+\chi(c'_j-a'_i)$, of the vertex weights
$w_A(a'_i)$ and $w_C(c'_j)$, with $w_A(y)=(1-\chi(b_0-y))(1-\chi(b_1-y))$ and
$w_C(y)=(1+\chi(b_0-y))(1+\chi(b_1-y))$, and of the indicator that the four labels are
distinct. Each weight has mean $1-q$, $L^2$ norm at most $2$ and supremum at most $4$;
testing a vertex by $\chi(z-y)$, for $z\notin\{b_0,b_1\}$, gives a weight of mean at most
$2r+2q$ in modulus (two quadratic character sums
contribute $q$ each and the cubic one $2r$),
with the same norm bounds. Expand the four internal edge factors into the $16$ subsets of
edges. A single edge, a Paley operator of norm $r$ between two weights of norm $2$, costs
$4r$; each of the two matchings costs $16q$; each of the four adjacent two-edge paths
costs $16q$ (integrate the two leaf Paley transforms first and use the supremum of the
internal weight); each of the four three-edge paths costs $64qr$; and the cycle costs
$16q$, as a product of two operators $M_{w_1}PM_{w_2}P$, each of Hilbert--Schmidt norm at
most $4r$. The distinctness costs only the same-side coincidences:
$w_A(y)w_C(y)=(1-\chi(b_0-y)^2)(1-\chi(b_1-y)^2)=0$ kills the four cross-side ones, and on a
same-side coincidence $w^2\le4w$ and every doubled edge factor is at most twice the
original, so its mass is at most
$16q\,\E_a[w_A(a)(\E_cw_C(c)(1+\chi(c-a)))^2]\le16q(1+8r+16q)$; the two such events cost
$C_{\rm coll}=32q+256qr+512q^2$. Without tests this bounds the normalizing mass below by
$(1-q)^4-16r-112q-256qr-C_{\rm coll}=\Delta_{\rm core}$. With the two vertices of $c'$
tested, the empty graph gives $m^2$, the four single edges $16rm$, the matchings $32q$, the
two paths centred on the untested side $32q$, the two paths centred at a tested vertex
$32qm$, the three-edge paths $256qr$, the cycle $16q$ and the coincidences $C_{\rm coll}$; the
sum is
\[
 m^2+16rm+80q+32qm+256qr+C_{\rm coll}=148q+616qr+580q^2 .
\]
Hence $|\E[\psi_{c'}(z)\mid b]|\le(148q+616qr+580q^2)/\Delta_{\rm core}$ for
$z\notin\{b_0,b_1\}$, and the two exceptional labels cost $2q$, which gives the soft bound
$c_{\rm core}$. The lower bound $\Delta_{\rm core}$ increases with $p$ and is positive at
$p=2^{20}$, hence for $p\ge2^{20}$.

\emph{The private score.} $\E[\Hs1\mid x]=r^{-1}\E_z\psi_b(z)\E[\psi_{b'}(z)\mid C]$, and
given $C$ the pair $b'$ is a uniform ordered distinct pair from the set $L_C$ of cyclic
extensions of $C$, of $n=p\bar{\mathsf K}(C)/16\ge p(1-\eta)/16$ points. For $z\notin C$,
expanding $\mathsf K(t,C)\chi(z-t)$ in $t$ gives $|\E_{t\in L_C}\chi(z-t)|\le\varphi_b$: the four
terms with one further character cost $q$ each, the six cubic terms $2r$ each, the four
quartic terms $2r+q$ each (\cref{thm:weil}(b)) and the quintic term $4r$, and the mass is
$\bar{\mathsf K}\ge1-\eta$. Sampling without replacement gives
$\E\chi(z-b'_0)\chi(z-b'_1)=(n\gamma_C^2-\E_{t\in L_C}\chi(z-t)^2)/(n-1)$ with
$\gamma_C=\E_{t\in L_C}\chi(z-t)$, of modulus at most $(1+d_b)\varphi_b^2+d_b$, as
$1/(n-1)\le d_b$; the four labels $z\in C$ cost $4q$. This gives the soft bound $c_b$, and
$c_{\rm avg}$ for $\Hbar$.

\emph{The row of the actual kernel.} (``Physical'' in this proof
means: of the actual kernel $K_0$, as opposed to the soft kernel
$K_{\rm soft}$.) The soft row of $K_{\rm soft}\circ\Hbar$ at $x$ is at most
$c_{\rm avg}\,\mathfrak r(C)\E\mathfrak r\le\mathfrak r_+\bar{\mathfrak r}\,c_{\rm avg}$ in modulus, and
the soft row of $K_{\rm soft}$ is at least $\mathfrak r_-^2$. Passing to the physical
conditional law removes, first, the eight coincidences of $b'$ with the core $C'$, which
are excluded from the actual half: given $C$ the pair $b'$ is uniform in $L_C$, so each
costs at most $16q/(1-\eta)$ of the row mass $\mathfrak r(C)\E\mathfrak r$, together
$I_{\rm in}$; second, the coincidences across the cut, of row mass at most $D_{\rm bal}$
(\cref{prop:l20-collisions}(a), the compression factor $h_p\le1$ being used in the
conservative direction); third, the equal-pair spikes, of score at most $1/r$ and row mass
at most $\mathfrak s$. On the first two kinds $|\Hbar|\le2+r$. The ratio of the remaining
numerator and denominator is at most $c_{\rm cond}$.
\end{proof}

The conditional bound $c_{\rm cond}$ is about $299r$; it is used only for the range of the
tilted marginal, by Jensen's inequality in \cref{lem:tilt-marginals}, while the much
smaller $L^2$ bound $\mathfrak l$ prices its variance.

\begin{proposition}[The chamber cut of a positive root]\label{prop:root-chamber}
Let $p\ge2^{30}$, suppose $k_4r+N_A<1$ and $\rho_{\rm in}>0$, and let
$\Gamma_{p,\ulaw}$ be the individual-square root, with
$s=1+\ulaw^2$, $\lambda=2\ulaw/s$, and $g=g(\lambda)$ of
\eqref{eq:root-cubic}. Put
\begin{gather*}
 \eps_0=\frac{2rd_{\mathfrak r}(\sqrt{30}\sqrt{g_{\mathbf z}}+2)+k_4^2q+2\mathfrak r_2N_A
   +2k_4rN_A+N_A^2}{h'_p}+D_{\rm bal},\qquad
 \eps_\ulaw=s\eps_0+2\ulaw\eps^{\rm raw}_{\rm lin}+\ulaw^2\eps^{\rm raw}_{\rm sq},\\
 z_-=s-2\ulaw h_0-\ulaw^2\eps_{\rm sq},\qquad z_{\rm raw}=(1-\eps_Z)z_-,\qquad
 g_{\rm ch}=\frac{sg\,r\Gamma_{20}+\eps_\ulaw}{z_{\rm raw}},\\
 c_{\rm ch}=\frac{s(\sqrt{30}+\lambda/3)r\Gamma_{20}+\eps_\ulaw+sd_{\mathfrak r}^2/h'_p}{z_{\rm raw}},
 \qquad
 d_\ulaw=\frac{s\sqrt{\chi^+_p}+2\ulaw\mathfrak l+\ulaw^2\eps_{\rm sq}}{z_-}.
\end{gather*}
If $z_->0$, $\ulaw c_{\rm cond}<1$, $c_{\rm ch}<1$ and $d_\ulaw<1$, the cut operator
$T_\ulaw$ of $\Gamma_{p,\ulaw}$ at $\vartheta_0$ in the fixed uniform actual half chart
$L^2(\mu_p)$ satisfies
\begin{gather*}
 -\Bigl(g_{\rm ch}+\frac{d_\ulaw^2}{1+g_{\rm ch}}\Bigr)I\preceq T_\ulaw\preceq
 \Bigl(1+\frac{d_\ulaw^2}{1-c_{\rm ch}}\Bigr)I,\qquad \|T_\ulaw\one-\one\|^2\le d_\ulaw^2,\\
 \frac{\sup T_\ulaw\one}{\inf T_\ulaw\one}\le R_p\,\frac{(1+\ulaw(2+r))^2}{(1-\ulaw c_{\rm cond})^2}.
\end{gather*}
With $g^\ulaw_{20}=(g_{\rm ch}+d_\ulaw^2/(1+g_{\rm ch}))/r$ and
$c^{+,\ulaw}_{20}=d_\ulaw^2/(q(1-c_{\rm ch}))$ evaluated at an endpoint $P$, the same bounds,
$-g^\ulaw_{20}p^{-1/2}\le T_\ulaw\le1+c^{+,\ulaw}_{20}/p$, the variance $d_\ulaw^2$ and the
range ratio hold for every prime $p\ge P$ with the same $\ulaw$. The block caps follow from
\cref{lem:spec:tensorcomp}(b): a middle block $\Blk$ of a piece starting at $X\ge P$, with
least prime $b$ and $\sum_{p\in\Blk}\log p\le L_{\rm loc}$, has
$T_{\Blk,\ulaw}\succeq-g^\ulaw_{20}\exp\bigl(c^{+,\ulaw}_{20}L_{\rm loc}/(X\log X)\bigr)b^{-1/2}$
if $g^\ulaw_{20}X^{-1/2}<1$.
\end{proposition}

\begin{proof}
\emph{The base kernel.} Expand $K_{\rm soft}=T[\mathsf R,\mathsf R]$ by \eqref{eq:l20-rectangle}
and keep the positive rank: $T[\mathfrak r\otimes\one,\mathfrak r\otimes\one]$ is the rank-one
operator $(\mathfrak r\otimes\one)\ip{\mathfrak r\otimes\one}{\cdot}$. The two rank/star classes
are the star part $T_\star=r\{\sqrt{30}(\mathbf v\mathbf z^*+\mathbf z\mathbf v^*)
+2(\mathbf w\mathbf e^*+\mathbf e\mathbf w^*)\}$ of \eqref{eq:l20-rank4-operator}, up to the
factor $\mathfrak r$ on one side, which costs $2rd_{\mathfrak r}(\sqrt{30}\sqrt{g_{\mathbf z}}+2)$;
star/star costs $k_4^2q$, rank/remainder $2\mathfrak r_2N_A$, star/remainder $2k_4rN_A$ and
remainder/remainder $N_A^2$, as $\|T[X,Y]\|\le\|X\|\|Y\|$. The compression to $\mu_p$
(\cref{lem:restriction}) and the collisions give, in the coefficient space,
\[
 P_0K_0P_0=r\,\mathcal FM\mathcal F^*+P_0\mathfrak r_\mu\mathfrak r_\mu^*P_0+\mathcal R_0,
 \qquad\|\mathcal R_0\|\le\eps_0,
\]
where $\mathfrak r_\mu$ is the function $\mathfrak r(C)$ on the actual half; the centred rank
$P_0\mathfrak r_\mu\mathfrak r_\mu^*P_0$ is positive semidefinite, of norm at most
$\|\mathfrak r-1\|_2^2/h_p\le d_{\mathfrak r}^2/h'_p$, and enters only the absolute norm.

\emph{The tilt.} By \cref{def:positive-root},
$K_0\circ w_\ulaw=K_0+2\ulaw\,K_0\circ\Hbar+\ulaw^2\,K_0\circ V$
with $V=\frac13\sum_j\Hs j^2$. By \cref{lem:root-normalization}(a),
$K_0\circ V=K_0+\mathcal R_V$ with
$\|\mathcal R_V\|\le\eps^{\rm raw}_{\rm sq}$, and $P_0(K_0\circ\Hbar)P_0$ is the pair
term up to $\eps^{\rm raw}_{\rm lin}$. As $1+\ulaw^2=s$ and $2\ulaw/s=\lambda$,
\[
 P_0(K_0\circ w_\ulaw)P_0=s\,r\,\mathcal FM_\lambda\mathcal F^*
 +sP_0\mathfrak r_\mu\mathfrak r_\mu^*P_0+\mathcal R_\ulaw,\qquad\|\mathcal R_\ulaw\|\le\eps_\ulaw .
\]
By \cref{lem:root-pair-matrix}, applicable since $\lambda\le1$,
$M_\lambda\succeq-gI$ and $\|M_\lambda\|\le\sqrt{30}+\lambda/3$, so
$\mathcal FM_\lambda\mathcal F^*\succeq-g\Gamma_{20}$ and
$\|\mathcal FM_\lambda\mathcal F^*\|\le(\sqrt{30}+\lambda/3)\Gamma_{20}$, the features mapping into
mean-zero functions.

\emph{Mass, marginal and range.} Apply \cref{lem:tilt-marginals} to the kernel of $T_0$,
which has mass one, $\|T_0\one-\one\|\le(\chi^+_p)^{1/2}$
(\cref{lem:record:root-marginal}(ii)) and range ratio $\sup T_0\one/\inf T_0\one\le R_p$: by
\eqref{eq:record:root-count} the density is a ratio of completion counts between two
bounds whose quotient is $R_p$. Take the scores $\eta_j=\Hs j$, which the transposition
permutes (\cref{lem:positive-root-properties}(c)), and the constants
\[
 (c_V,\delta_0,e,\ell,h,c,b)=\bigl(1,(\chi^+_p)^{1/2},\eps_{\rm sq},\mathfrak l,
 h_0,c_{\rm cond},2+r\bigr)
\]
of \cref{lem:root-normalization} and
\cref{lem:positive-root-properties}(a). It gives $\E_{\Gamma_{p,0}}w_\ulaw\ge z_-$,
$\|T_\ulaw\one-\one\|\le d_\ulaw$ and the range ratio. Since the raw mass of
$K_0\circ w_\ulaw$ is $Z_0\,\E_{\Gamma_{p,0}}w_\ulaw\ge z_{\rm raw}$,
$P_0T_\ulaw P_0\succeq-g_{\rm ch}P_0$ and $\|P_0T_\ulaw P_0\|\le c_{\rm ch}$.

\emph{The constant block.} $T_\ulaw$ is self-adjoint with $\ip\one{T_\ulaw\one}=1$ and
$\|P_0T_\ulaw\one\|=\|T_\ulaw\one-\one\|\le d_\ulaw$, so \cref{prop:matching:positive}(a)
and (a$'$) give the two full bounds; no alignment of spectral directions is assumed.

\emph{Uniformity.} The quantities $\mathfrak r_-$, $h'_p$, the square term $4q\mathfrak r_2^2$
and the conditional ratio were coarsened so that every normalized remainder has a positive
increasing numerator and a positive decreasing denominator in $q,r$; the retained
$N_A/q$, $D_{\rm bal}/q$, frame cap, normalizer error, $\chi^+_p$ and $R_p$ have this
property (\cref{prop:l20-collisions}, \cref{lem:record:root-marginal}). Thus
$G=g_{\rm ch}/r$, $c_{\rm ch}/r$ and $D=d_\ulaw/r$ increase toward a smaller prime
endpoint. The negative coefficient is $G+D^2r/(1+Gr)$; its derivative in $G$ is
$1-D^2r^2/(1+Gr)^2>0$, as $Dr<1$, and its derivatives in $D$ and $r$ are nonnegative
(\cref{lem:majorant}(ii)). The
positive coefficient $D^2/(1-(c_{\rm ch}/r)\,r)$, the variance divided by $q$ and the range multiplier
increase as well on the gated domain, so the values at $P$ bound every prime $p\ge P$.
For the block caps, the primes of $\Blk$ satisfy the prime bounds and
$\sum_{p\in\Blk}1/p\le L_{\rm loc}/(X\log X)$ (\cref{lem:record:local-domain}), the cut
operator of $\Gamma_{\Blk,\ulaw}$ is the tensor product of the prime operators
(\cref{lem:spec:cutkernel}(f)), and \cref{lem:spec:tensorcomp}(b) applies.
\end{proof}

\subsection{The layer-flip symmetry \texorpdfstring{$\iota$}{iota} for the tilted law}\label{subsec:positive-root-layer}

This subsection bounds the reserved cut $\iota$ of a positive root, the cut L21. Here the
scores are positive Gram kernels, so the tilt costs only at order $r$; the bounds are
proved on the fixed uniform induced layer half, the half law of the one-layer root, which
is not tilted.

\paragraph{The two components.}
At $\iota$ the half graph of $G_{\rm L,2}$ is $G_{\rm L,1}$, two directed triangles on the
bases $\{0,2,4\}$ and $\{1,5,3\}$ (\cref{lem:layer12-geometry}(c)). Let $\mu_\triangle$ be the
uniform law of a directed cyclic triangle $x_0\to x_1\to x_2\to x_0$ of $\Fp$, let
$\nu_\triangle=\mu_\triangle\otimes\mu_\triangle$ (one triangle for each component), and let
$\nu=\lambda_p(G_{\rm L,1})$ be $\nu_\triangle$ conditioned on its six labels being distinct,
an event of probability $h\ge h_-=1-9q$. In a component, with the labels $x$ of the layer
$0$ and $x'$ of the layer $1$ indexed along the triangle, the cross edges are
$x_i\to x'_{i+1}$ and $x'_i\to x_{i+1}$: a six-cycle without an edge between the two labels
of one base. The soft kernel of a component is
\[
 \Phi(x,x')=\prod_{i\ne j}\bigl(1+(D_\triangle)_{ij}\chi(x'_j-x_i)\bigr),\qquad
 N_\Phi=\Phi-1,
\]
$D_\triangle$ being the skew adjacency matrix of the directed triangle
($(D_\triangle)_{i,i+1}=1$, $(D_\triangle)_{i+1,i}=-1$ with indices mod $3$, and
$0$ on the diagonal), and the soft layer
kernel on $\nu_\triangle\otimes\nu_\triangle$ is $\Phi\otimes\Phi$. A score $\Hs j$ multiplies the
kernel by $r^{-1}\ip{f_j(x)}{f_j(x')}$, $f_j(x)(z)=\chi(z-x_j)\chi(z-x_{j+3})$, the product of
the raw pair rows of the two selected labels of a half, one from each component. Put
\begin{gather*}
 d_2=\frac2{1-q},\quad d_\mu=\frac8{1-q^2},\quad
 \kappa=\frac{\sqrt3+\sqrt7}2,\\
 \mathfrak b=6+15d_2+18\sqrt{d_2d_\mu}+4d_\mu+14d_\mu r,\quad
 \mathfrak n=\kappa(1+5r)r+\mathfrak bq,\quad M_\Phi=d_\mu(1+q)^3,\\
 Z_-=(1-\mathfrak bq)^2-q\bigl[18M_\Phi^2+384M_\Phi+72d_\mu^2(1+q)^5\bigr],
\end{gather*}
and let $p\ge2^{30}$. The following facts hold for every score.
\begin{enumerate}[label=(L\arabic*)]
\item $\|N_\Phi\|\le\mathfrak n$ and $\|N_\Phi\one\|\le\mathfrak bq$ on $L^2(\mu_\triangle)$.
 Indeed the six single edges and the six two-edge stars have the coefficient matrix
 $\left(\begin{smallmatrix}iD_\triangle&-I_3\\-I_3&0\end{smallmatrix}\right)$
 (\cref{cor:star-matrix}), whose characteristic polynomial is $(z^2-1)(z^4-5z^2+1)$ and whose
 norm is $\kappa$. The six features (the coordinate isometries $V_i$ and the pair stars
 $W_{\{j,k\}}/r$) have actual Gram at most $(1+5r)I$ for $p\ge2^{20}$
 (\cref{lem:general-star}): the off-diagonal norms are at most $r\sqrt{1+q}/(1-q)$
 (singleton/singleton), $2q/(1-q)$ (incident singleton/pair), exactly $r$ (opposite
 singleton/pair: condition on that coordinate and negate while swapping the other two)
 and $(3q+\sqrt3r)/(1-q)$ (pair/pair), with row sums at most $5r$. The uncentred stars
 leave $+6q$ times the constant kernel, each pair mean being $-q$ and each star
 coefficient $-1$. The remaining $51$ masks can be priced using just
 their active coordinates. The one-, two-, and three-coordinate densities
 of a directed triangle are $1,d_2,d_\mu$. A kernel using $u$ output and
 $v$ input coordinates therefore incurs only $\sqrt{d_ud_v}$.
 The exact mask partition is
 \[
 \begin{array}{c|c|c|c}
 \text{mask type}&(u,v)&\text{number}&\text{bound per mask}\\ \hline
 \text{two-edge matching or forest}&(2,2)&15&d_2q\\
 \text{matching number two forest}&(2,3),(3,2)&9+9&\sqrt{d_2d_\mu}q\\
 \text{two disjoint two-edge stars}&(3,3)&3&d_\mu q\\
 \text{matching number three forest}&(3,3)&14&d_\mu rq\\
 \text{full six-cycle}&(3,3)&1&d_\mu q
 \end{array}
 \]
 A path of $k\le5$ edges has full-cube norm at most
 $r^{\lceil k/2\rceil}$: the two-edge star has squared norm
 $q(1-q)$, the three- and five-edge paths factor into Paley
 transforms with intervening contractions, and the four-edge path
 is bounded by fixing its middle vertex and applying two transforms.
 Disconnected components tensor. For the full cycle, fix one
 coordinate on each side and apply the two-edge matching/path
 bound, at cost $q$; no $r^3$ estimate for that cycle is used.
 Summing the table and the six means gives $\mathfrak b$.
 In particular
 \[
 \mathfrak b\le6+\frac{102}{1-q}
                     +\frac{32+112r}{1-q^2}<140.004
                     \qquad(p\ge2^{30}),
 \]
 with limit $140$. All these majorants are increasing in $q,r$.
\item The passive selected-coordinate maps $h\mapsto\E_z\chi(z-y)h(z)$ and
 $F\mapsto\E_{z,z'}\chi(z-y)\chi(z'-y)F(z,z')$ have norms $r$ and $(q(1-q))^{1/2}$
 (\cref{cor:raw-pair-maps}(a)).
\item \emph{Insertion of a score multiplier in a monomial of $N_\Phi\otimes N_\Phi$.}
In a monomial of $N_\Phi\otimes N_\Phi$ choose one cross
 edge in each component; it has at most one endpoint among the selected labels of a score,
 as no cross edge joins the two labels of a base. With the other vertices pinned, a bounded
 multiplier depending on at most one endpoint of each chosen edge can be placed between
 the two Paley transforms, or before or after their tensor product, so the kernel has norm
 at most $q$ times the supremum of the multiplier (\cref{lem:matching:kernel}(d)). Given the
 other two vertices of a triangle, the active vertex has exactly $(p+1)/4$ completions
 (\cref{lem:paley-triple}(a)), a conditional density at most $4/(1+q)\le4$; Cauchy--Schwarz
 over the pinned vertices recovers the triangle energies, at the cost $16$.
\item \emph{Rows and coincidences.} Each two-edge column of $\Phi$ has uniform mean $1+q$, so
 the rows of $\Phi$ are at most $M_\Phi$. A cross-component equality across the cut costs at
 most $4d_\mu^2(1+q)^5q$ of row mass (given a coordinate, the other two of its triangle have
 density $d_\mu$), a twin equality $x_u=x'_u$ at most $64M_\Phi q$, an equality between
 the two components inside one half at most $M_\Phi^2q$ (integrate the other half against
 the two row bounds), and a same-component equality at different bases makes $\Phi$ vanish,
 the third base supplying a contradictory edge. The soft mass is at least
 $(1-\mathfrak bq)^2$ by (L1); counting the eighteen cross-component coincidences inside the
 two halves, the six twins and the eighteen cross-component coincidences across the cut,
 the raw mass of the uniform root relative to $\nu\otimes\nu$ is at least $Z_-$; division by
 the half probability squared only increases it.
\end{enumerate}

\begin{lemma}[Inserting an opposite-base-pair score into the soft kernel]\label{lem:native-pair-linear}
Fix one of the opposite-base-pair scores $\Hs0,\Hs1,\Hs2$ introduced
above and denote it by $\mathsf H$. On the two independent cyclic
triangles put
$K=\Phi\otimes\Phi$, $e=\mathfrak b-6$. Then, for $p\ge2^{30}$,
\[
 \|K\circ\mathsf H-\mathsf H\|\le C_{\rm pair}q,\qquad
 C_{\rm pair}=2304+2\kappa(1+5r)(1+e)
                    +r(12+14e+e^2)<2897.
\]
Restriction to the actual six-distinct half costs at most
$(1-9q)^{-1}$. Cross-cut collisions must still be charged separately.
\end{lemma}
\begin{proof}
Keep the same auxiliary label in the two raw pair rows. For an
operator $A$ on the first triangle, let
\[
 (VF)(x,y)=\E_s\chi(s-x_i)\chi(s-y_j)F(x,s).
\]
Uniformity of $y_j$ and the Paley bound give $\|V\|\le r$, while
direct composition gives
$(A\otimes J)\circ\mathsf H=r^{-1}V(A\otimes I)V^*$.
Thus this term has norm at most $r\|A\|$.
Write $\Phi=J+S+R$, with $S$ the twelve raw star monomials.
The preceding bounds give
$\|S\|\le\kappa(1+5r)r+6q$ and $\|R\|\le eq$.
Each nonempty old-edge monomial times $\mathsf H$ costs at most
$16q$: select one old edge and the passive score coordinate in
the other triangle. The passive Paley map and the two-edge
matching/path bound leave $r^3$ before the raw factor $r^{-1}$;
the conditional triangle densities cost $16$. The absence of an
old edge at the selected base is used here.
Apply this bound to the $12^2=144$ terms of $S\otimes S$.
The one-sided $S$ and $R$ terms use the displayed sandwich.
The remaining $R\otimes S$, $S\otimes R$, $R\otimes R$ terms
use the Schur multiplier norm at most $r^{-1}$ of the positive
raw pair Gram. Summing gives the stated bound $C_{\rm pair}q$.
The endpoint inequality uses $r\le2^{-15}$ and the rational
upper bound for $\mathfrak b$ above. Conditioning both half laws
on six distinct labels gives the single factor $(1-9q)^{-1}$.
\end{proof}

\begin{proposition}[The layer cut of a positive root]\label{prop:root-layer}
Let $p\ge2^{30}$ and let $\Gamma_{p,\ulaw}$ be the individual-square
root, with $\gamma=\ulaw^2$. Let $\nu$ be the fixed uniform induced layer
half. With $d_\mu,\mathfrak b,\mathfrak n,M_\Phi,Z_-$ as above and
$h_0$ of \cref{lem:root-normalization}, put
$B_H=2+r$ and
\begin{gather*}
 m_\ulaw=1+2\ulaw r+\gamma,\qquad F_\ulaw=1+2\ulaw B_H+\gamma B_H^2,\qquad
 C_{\rm nn,\ulaw}=3969\cdot16\,F_\ulaw,\\
 C_{\rm sc,\ulaw}=127\bigl[4(1+4\ulaw r+4\gamma q)+2(q+2\ulaw r+\gamma)\bigr],\qquad
 C_{\rm cr,\ulaw}=\bigl[256M_\Phi+64d_\mu^2(1+q)^5\bigr]F_\ulaw,\\
 C_\ulaw=C_{\rm nn,\ulaw}+C_{\rm sc,\ulaw}+C_{\rm cr,\ulaw},\qquad
 C_{\rm in,\ulaw}=127\bigl[9F_\ulaw+2(q+2\ulaw r+\gamma)\bigr],\\
 z^\iota_\ulaw=1+\gamma-\gamma\Bigl[4q+36qB_H^2+\frac{2q}{1-q}\Bigr]-2\ulaw h_0,
\end{gather*}
and
\begin{align*}
 g^\iota_\ulaw&=\frac{2\mathfrak nm_\ulaw+C_\ulaw q}{rz^\iota_\ulaw h_-Z_-},\\
 c^\iota_\ulaw&=\frac{2\ulaw r+\gamma c_0q
   +(2\mathfrak nm_\ulaw+C_\ulaw q)/h_-}{z^\iota_\ulaw Z_-},\\
 A^\iota_\ulaw&=\frac{C_\ulaw+C_{\rm in,\ulaw}+910\gamma
   +2(1+\gamma)\mathfrak b/h_-}{z^\iota_\ulaw h_-Z_-},\qquad
 c^{+,\iota}_\ulaw=\frac{q(A^\iota_\ulaw)^2}{1-c^\iota_\ulaw}.
\end{align*}
Here $C_{\rm nn,\ulaw}$ bounds the $N_\Phi\otimes N_\Phi$ monomials,
$C_{\rm sc,\ulaw}$ the selected-label coincidences, $C_{\rm cr,\ulaw}$ the other
cross coincidences, and $C_{\rm in,\ulaw}$ the input-half coincidences. If
$z^\iota_\ulaw>0$ and $c^\iota_\ulaw<1$, the cut operator $T^\iota_\ulaw$ of
$\Gamma_{p,\ulaw}$ at $\iota$ on $L^2(\nu)$ satisfies
\[
 -g^\iota_\ulaw rI\preceq T^\iota_\ulaw\preceq(1+c^{+,\iota}_\ulaw q)I,\qquad
 \|P_0T^\iota_\ulaw\one\|\le qA^\iota_\ulaw,\qquad\|P_0T^\iota_\ulaw P_0\|\le c^\iota_\ulaw .
\]
The bounds at an endpoint hold for all larger primes with the same $\ulaw$, and
$g^\iota_\ulaw\to\sqrt3+\sqrt7$. The block caps follow from \cref{lem:spec:tensorcomp}(b), as in
\cref{prop:root-chamber}.
\end{proposition}

\begin{proof}
Write $w_\ulaw=\frac13\sum_j\Upsilon_j$, with
$\Upsilon_j=1+2\ulaw\Hs j+\gamma\Hs j^2$. It suffices to bound the soft
operator of one polynomial $\Upsilon=1+2\ulaw\mathsf H+\gamma\mathsf H^2$
and then average and normalize. By \cref{lem:spec:cutkernel}(f),
$T^\iota_\ulaw=(T^\iota_0\circ w_\ulaw)/\E_{\Gamma_{p,0}}w_\ulaw$.

\emph{Soft expansion.} $\Phi\otimes\Phi=1\otimes1+N_\Phi\otimes1+1\otimes N_\Phi+N_\Phi\otimes N_\Phi$.
The base term $(1\otimes1)\circ\Upsilon$ is positive semidefinite: $\mathsf H$ is a positive Gram kernel,
so is its pointwise square, and the coefficients of $\Upsilon$ are nonnegative. In
$(N_\Phi\otimes1)\circ\Upsilon$ the passive component carries the constant kernel and its selected
labels; retaining the auxiliary variables, (L1) and (L2) give the absolute price
$\mathfrak n[1+2\ulaw r+\gamma(1-q)]\le\mathfrak nm_\ulaw$ in both orders, and the other
component the same. For the $63^2$ monomials of $N_\Phi\otimes N_\Phi$, (L3) with the
multiplier $\Upsilon$ on four distinct selected labels, where $|\mathsf H|\le B_H$ (\cref{thm:weil}(b)) and
$|\Upsilon|\le F_\ulaw$, gives $16F_\ulaw q$ each, together $C_{\rm nn,\ulaw}q$. The strict score
mask is part of this multiplier; no other cross-injectivity mask is removed inside a norm.

\emph{Coincidences.} The four single coincidences between the selected labels of the two
halves have, before the input half is conditioned, total row and column mass at most $4q$,
and there $|\mathsf H|\le2r$, so $|\Upsilon|\le1+4\ulaw r+4\gamma q$; the two equal-pair events have mass
at most $2q^2$ and $\mathsf H=(1-2q)/r$, so $|\Upsilon|\le1+2\ulaw/r+\gamma/q$. They are deleted from the
unmasked part $1+N_\Phi\otimes1+1\otimes N_\Phi$, whose kernel is pointwise
at most $1+2\cdot63=127$ in absolute value ($N_\Phi$ being a sum of $63$
character monomials of modulus at most $1$), which costs $C_{\rm sc,\ulaw}q$. After these four, the coincidences across the cut are four twins and
sixteen cross-component equalities, on which $|\Upsilon|\le F_\ulaw$, and (L4) gives
$C_{\rm cr,\ulaw}q$. These are positive row and column bounds of the deleted kernels, not
Schur monotonicity of signed kernels.

\emph{Mass.} The raw mass of the uniform root is at least $Z_-$ (L4). Its scores have
$|\E_{\Gamma_{p,0}}\Hs j|\le h_0$ and $\E_{\Gamma_{p,0}}\Hs j^2\ge1-4q-36qB_H^2-2q/(1-q)$
(\cref{lem:root-normalization}(b)); so $\E_{\Gamma_{p,0}}w_\ulaw\ge z^\iota_\ulaw$.

\emph{Lower bound.} The restriction to $\nu$ costs $h_-^{-1}$ (\cref{lem:restriction}) and the
normalization $Z_-^{-1}$ of the uniform root. Dividing by the tilt mass gives $T^\iota_\ulaw\succeq-g^\iota_\ulaw r$.

\emph{Centred norm.} The selected pair of a half is a uniform distinct pair under $\nu$, by
square-affine invariance and the interchange of the two components, so the centred base
is at most $2\ulaw r+\gamma c_0q$: $P_0\mathsf HP_0=r\,\mathcal U\mathcal U^*$ with
$\|\mathcal U\|\le1$, and \cref{cor:raw-pair-maps}(b) gives $\|P_0\mathsf H^2P_0\|\le c_0q$. The other
terms are those of the lower bound, which gives $c^\iota_\ulaw$.

\emph{Row.} On a six-distinct output half the unconditioned base row is constant, as
$\mathsf H$ has mean $0$ and $\mathsf H^2$ has mean $(1-2q)^2+q^3$ when the other half is
unconditioned. In a one-component term the linear row vanishes
after the passive coordinate is integrated, and the square row multiplies $N_\Phi(x,x')$ by
\[
 \E_z\bigl[\chi(z-y)^2\chi(z-y')^2\chi(z-y'')^2\bigr]-\bigl(\E_z\chi(z-y)\chi(z-y')\chi(z-y'')\bigr)^2,
 \qquad y\ne y'',
\]
where $y,y''$ are the two selected labels of the output half and $y'$ that
of the input half in the active component; this differs from one by at
most $7q$, also when $y'\in\{y,y''\}$. With $|N_\Phi|\le65$ this
costs $455q$. The two
components therefore have row price $2(1+\gamma)\mathfrak bq/\sqrt h+910\gamma q$. The nine
input-half coincidences and their pair spikes cost $C_{\rm in,\ulaw}q$, by the same
ordinary/equal-pair split; all other errors cost $C_\ulaw q$. With $h^{-1/2}\le h_-^{-1}$, the
half compression, the normalization of the uniform root and the tilt
mass this gives $\|P_0T^\iota_\ulaw\one\|\le qA^\iota_\ulaw$.

\emph{Upper bound.} $T^\iota_\ulaw$ is self-adjoint of mass one, so
\cref{prop:matching:positive}(a) gives $T^\iota_\ulaw\preceq1+q^2(A^\iota_\ulaw)^2/(1-c^\iota_\ulaw)$.

\emph{Uniformity.} After the explicit factor $r$ in $g^\iota_\ulaw$ is cancelled, all numerators
are positive increasing in $q,r$ and all lower denominators positive decreasing, so
evaluation at $q=1/P$, $r=P^{-1/2}$ is uniform for $p\ge P$. As $q\to0$, $\mathfrak n/r\to\kappa$
and $m_\ulaw$ and $z^\iota_\ulaw$ tend to $1+\gamma$, so $g^\iota_\ulaw\to2\kappa=\sqrt3+\sqrt7$.
\end{proof}

For $\gamma=1$, the individual squares at $\ulaw=1$, the proposition is
the individual-square layer bound at that parameter; then
$z^\iota_1=2-4q-36qB_H^2-2q/(1-q)-2h_0$, $2\mathfrak nm_1/r=2[\kappa(1+5r)+\mathfrak br](2+2r)$
and $F_1=(3+r)^2$. The proved coefficient tends to $2\kappa$ for every $\ulaw>0$, because the
two one-component terms are priced separately; the larger layer cap costs only through
the weight $v_{\rm L}$ of the completion (\cref{rem:middle-weights}).

\paragraph{Sharper bounds for the tilted law (next nine subsections).}
The following subsections sharpen the two cut bounds above at the
seven interval starts. In order they: isolate collision terms on
two explicit subspaces; give a finite gap for the uniform law;
split a remainder across that gap; retain a constant term at the
layer-flip cut; list supports for the $39$-trace law; extend it to
all sorting cuts; use one parameter for both cuts; evaluate an
explicit layer price; and evaluate a finite chamber price. A first reading can skip to
\cref{subsec:positive-root-caps} and use the stated bounds.

\subsection{Isolating collision terms on two explicit subspaces}
\label{imp:local:protected-section}

The remaining subsections of this section are finite calculations
at one prime $p\equiv3\pmod4$, used only to evaluate the caps
tabulated in \cref{subsec:positive-root-caps}. They do two things.
First, they sharpen the collision part of the L20 bound for the
uniform principal law, by computing it exactly on two explicit
subspaces. Second, they construct and bound the coherent tilted
law of the last interval $[2^{45},P_3)$: its score is an explicit
linear combination of the trace functions
$\E_s\prod_{i\in D}\chi(s-x_i)$ over a list of $39$ supports
$D$ (subsets of the twelve label positions), symmetrised so that
one law serves at every sorting cut. Its matrix bounds are first
proved at the endpoint parameter $t_*=12/5$ and then transferred
by convexity to the parameter $t=67/40$ used in the construction.
In this part of the paper a support or a law is called
\emph{physical} when it is described in the original twelve
labels, as opposed to coordinates of a chosen half chart; a
\emph{gate} is one finite matrix or positivity inequality checked
exactly (\cref{app:certificates}). A reader who accepts the caps of
\cref{cert:positive-root-caps} can skip to
\cref{sec:middle-frame}. Local words in this part: a \emph{pin}
is a label fixed by conditioning (core pins $C_i$, private pins
$b_l$); a \emph{mask} is a subset of the edge set of the small
graph under discussion, determining one character product; an
\emph{arm} is one synthesis factor (feature map) in a product of
two stars, and a \emph{free leaf} is a label coordinate occurring
in no other factor; a \emph{spike} is an exceptional large value
of a trace or score at special configurations (in particular at
equal score pairs), always bounded separately; a \emph{hole
term} is a term in which the common auxiliary label equals one of
the fixed labels; an argument is \emph{rooted} when it keeps the
common auxiliary summation label $s$ unintegrated, and
\emph{unrooted} when $s$ is integrated out first; a star factor
is \emph{passive} when it carries none of the displayed
characters; the \emph{realification} of a complex feature space
is the real span of the real and imaginary parts of its
functions; an \emph{endpoint} $X$ is the lower endpoint of an
interval, at which bounds are evaluated; and \emph{physical}
(as defined above) is opposed to \emph{soft}: the physical
kernel is the actual kernel $K_0$ in the original labels, the
soft kernel its expansion without distinctness.

We begin with the uniform law, before any tilt. The collision
part of its L20 kernel---the terms where labels of the two halves
coincide---is split according as the test function lies in one of
two explicit subspaces or in their orthogonal complement; on the
two subspaces it is computed exactly below.
In the following local calculations write $A=\mathsf R$, $a=\mathfrak r$
for the rectangle and its private mean in \cref{prop:l20-rectangular},
$C_0=A-a\otimes1$, and $Z=\boldsymbol z$, with the other normalized
star features $v,w,e$ as in \cref{prop:l20-rank4}. All spaces have
probability normalization, $q=p^{-1}$ and $r=\sqrt q$.
The letter $a$ in these local formulas denotes a function on the core;
the scalar exponent does not enter this calculation. Likewise $t$
below denotes the root parameter $\ulaw$, not scalar repricing.
Write $Q=I-J$ for ordinary centering (subtracting the mean). Put $a_2=\|a\|_2$, $P_a=aa^*/a_2^2$ and
\[
 P_{\rm c}=I_{\rm core}\otimes J_{\rm pair},\qquad
 P_{\rm p}=P_a\otimes Q_{\rm pair},\qquad
 P=P_{\rm c}+P_{\rm p},\qquad R=Q_a\otimes Q_{\rm pair}.
\]
The core labels are $C=(a_0,a_1,c_0,c_1)=(C_0,C_1,C_2,C_3)$, with $a_i\to c_j$,
and the private pair is $b=(b_0,b_1)$. Set
\[
 K_C(s)=\prod_i(1-\chi(s-a_i))\prod_j(1+\chi(s-c_j)),
 \qquad\kappa_C=\E_sK_C(s),\qquad
 Z_C=(K_C-\kappa_C)/(\sqrt{15}r).
\]
The six-distinct induced half is obtained by conditioning the
independent core and distinct private pair on the six labels being pairwise distinct.

\begin{lemma}[Compression of the constant kernels at pinned labels]
\label{imp:local:pin-projections}
For $h\in H=L^2_0(\mathbb F_p)$, let $\sigma=(-1,-1,1,1)$, $V_i h=h(C_i)$,
$B_lh=h(b_l)$, $w=\sum_iV_i/2$, $u=(V_0+V_1-V_2-V_3)/2$,
$v=(B_0+B_1)/\sqrt2$, and
\[
 W h(b)=\mathbb E_s\chi(s-b_0)\chi(s-b_1)h(s),\quad e=W/r,
 \quad R_lh(b)=\mathbb E_s\chi(b_l-s)h(s).
\]
On $H$, for $\epsilon\in\{-1,1\}$, define
\[
 F_i^\epsilon h=\prod_{l=0}^1(1+\epsilon\chi(b_l-C_i))h(C_i),
 \qquad G_lh=K_C(b_l)h(b_l),
\]
and set $k=\mathbb E a/a_2$, $k'=\mathbb E(a\kappa)/a_2$,
where $\kappa$ denotes the function $C\mapsto\kappa_C$ and the
expectations are under the core law (uniform on core
configurations).
The sixteen core-pin constant kernels and the four private-pin
constant kernels have auxiliary-mean-zero parts
\[
 K_c=q\sum_{i,j}F_i^{\sigma_j}(F_j^{\sigma_i})^*,
 \qquad K_b=q(G_0+G_1)(G_0+G_1)^*.
\]
Their protected blocks are
\begin{equation}\label{eq:merge-local-I-14}
\begin{array}{c|c|c|c}
 &P_{\rm c}KP_{\rm c}&P_{\rm p}KP_{\rm p}&P_{\rm c}KP_{\rm p}\\ \hline
K_c&4q(1-q)^2ww^*&16q^2k^2ee^*
 &4q(1-q)k[-2u(R_0+R_1)^*+2wW^*]\\
K_b&60q^2ZZ^*&2q(k')^2vv^*&2\sqrt{30}qrk'Zv^*
\end{array}
\end{equation}
The full auxiliary constant-label pieces contribute, respectively,
$16q(1-q)^2J$ and $4q\kappa\otimes\kappa$ on $P_{\rm c}$, and zero
other protected blocks.
\end{lemma}

\begin{proof}
Uniformity of the private pair and affine invariance of the actual
core give exactly
\[
 P_{\rm c}F_i^\epsilon=(1-q)V_i,\quad
 P_{\rm p}F_i^\epsilon=k[\epsilon(R_0+R_1)+W],\quad
 P_{\rm c}G_l=\sqrt{15}rZ,\quad P_{\rm p}G_l=k'B_l.
\]
Indeed $\mathbb E[a\mid C_i=s]=\mathbb E a$ and
$\mathbb E_C[aK_C(s)]=\mathbb E(a\kappa)$, independently of $s$.
Summing uses $\sum\sigma_i=0$, $\sum V_i=2w$, and
$\sum\sigma_iV_i=-2u$, proving \eqref{eq:merge-local-I-14}.  On the constant label,
$P_{\rm c}F_i^\epsilon1=(1-q)1$, $P_{\rm p}F_i^\epsilon1=0$,
$P_{\rm c}G_l1=\kappa$, and $P_{\rm p}G_l1=0$, which proves the last assertion.
\end{proof}

\subsection{A finite collision bound for the uniform law}
\label{imp:local:finite-collision-section}

This result is a finite bound for the uniform (untilted) L20 root in the original twelve labels.
Its positive coefficient, marginal range, variance and layer data
are consequently unchanged. Retain the independent-half rectangle
$A=a\otimes1+C_0$ and the certified bounds for the rectangle
and its domain gap from \cref{prop:l20-rectangular} and
\cref{prop:l20-pair-gap}: numbers again called $a_2,c$ with
$a_2\ge\|a\|_2$, $c\ge\|C_0\|$, $z_2\ge\|Z\|^2$ (here $a_2,c$
are upper bounds, although the same letters denoted the exact
norms above).
Put
\[
 \theta=(1-q)(1-3q)(1-7q),\quad h_-=(1-4q)(1-5q),
 \quad \eta=11r+2q,
\]
\begin{equation}\label{eq:merge-local-I-15}
 \begin{aligned}
 d_\kappa&=\sqrt{\frac{5q+34qr+11q^3}{\theta}}+2q,\qquad
 d_a=\frac{(2+\eta+16q)d_\kappa+15q}{1-q},\\
 a_-&=(1-\eta)(1-\eta-16q),\qquad \omega=d_a/a_-.
 \end{aligned}
\end{equation}
From the rectangle estimates of \cref{sec:middle-reference} we
retain the bounds $|\kappa-1|\le\eta$, $\|a-1\|_2\le d_a$ and
$a\ge a_-$.

\begin{proposition}[Finite bound for the collision term on the two protected subspaces]
\label{imp:local:collision-gap}
Define
\[
 \delta_c=\frac{4(4r+92q)}{\sqrt{\theta(1-q)}},
 \qquad \delta_b=\frac{4\sqrt{15}\,r}{\sqrt{\theta(1-q)}},
\]
\begin{equation}\label{eq:merge-local-I-16}
 E_{\rm pin}=256q(2\delta_c+\delta_c^2)
              +1024q(2\delta_b+\delta_b^2),\qquad
 E_{\rm mult}=\frac{24903680q^2}{1-q},
\end{equation}
and
\begin{equation}\label{eq:merge-local-I-17}
 \begin{aligned}
 b_B&=64q(2+\eta),\\
 A_0&=4q+60q^2z_2,\qquad
 A_1=2q(1+\eta)^2+16q^2,\\
 X&=8qr\bigl(\sqrt{2/(1-q)}+1\bigr)
                +2\sqrt{30}qr(1+\eta)\sqrt{z_2},\\
 \alpha&=\max(A_0,A_1)+X+16q\omega^2+4q(\eta+\omega)^2,\\
 E_S&=\frac{b_B^2}{a_-^2-b_B},\qquad
 G=a_2c+\alpha,\qquad D=c^2+b_B .
 \end{aligned}
\end{equation}
For $p\equiv3\pmod4$ satisfying $a_-^2>b_B$ and $G>D$, the old mass-one root operator
on its fixed induced half satisfies
\begin{equation}\label{eq:merge-local-I-18}
 T_{\rm old}\succeq-g_{\rm pin}rI,\qquad
 g_{\rm pin}=
 \frac{E_S+G+b_B^2/(G-D)+E_{\rm pin}+E_{\rm mult}}
      {h_-(1-\varepsilon_Z)r}.
\end{equation}
Here $\varepsilon_Z$ is the normalizer error of the uniform
(untilted) law from \cref{prop:l20-completion}. Exact rational endpoint evaluation is valid uniformly
for every prime above the endpoint. For example, the endpoint $19951586$ gives
$g_{\rm pin}<5.481634404$, compared with the retained
$5.500057067029$.
\end{proposition}

\begin{proof}
Let $B$ be the sum of the twenty full-label constant pin-incidence
kernels in Lemma~\ref{imp:local:pin-projections}. For a core pin
$C_i=C'_j=t$, where $C'$ denotes the core labels of the second half, its incident factors are two half multipliers bounded
by four. After removing them, the cross polynomial is a swapped
product of two $3$-by-$2$ rectangles. Given one core coordinate,
the remaining three have density $16/\theta$ relative to the
three-label cube; the private distinct-pair density is $1/(1-q)$.

The six-edge rectangle has seventeen nonempty stars and
forty-six other masks (edge subsets that are not stars). Group the stars on the full independent cube.
Every raw star has norm at most $r$, because its Gram is a Schur
product of Paley Grams with diagonal at most one and norm at most
$q$. Distinct output supports are orthogonal ANOVA spaces.
On nonzero translation sectors, the three singleton and four
higher output supports, and the two singleton and one pair input
supports, give a seven-by-three matrix. Collapsing the two equal
singleton columns gives Gram
\[
 \begin{pmatrix}14&\pm i\sqrt2\\\mp i\sqrt2&3\end{pmatrix},
\]
whose largest eigenvalue is below sixteen: the three retained
core signs sum to $\pm1$. On zero translation frequency the
blocks are $2\sum_{|S|\ge2}\sigma_Sm_S\otimes1$, of norm at most
$4r$, and $3\,1\otimes m_{\rm pair}$, of norm at most $3r$.
Their input and output spaces are orthogonal. This retains the
full auxiliary constant sector and proves that the seventeen stars together have norm at most $4r$.
Every other mask contains two disjoint edges whose induced
subkernel is a matching, a path, or a restricted $K_{2,2}$, of
norm at most $2q$.
Fixing the third output coordinate inserts only one-variable input
multipliers and separate output exclusions. Distinctness of the
selected input and output pairs stays inside the restricted
kernel. Density congruence gives $\|A_i-J\|\le\delta_c$.

For a private pin the two incident factors are bounded by sixteen.
The remaining rectangle is $K_C(s)$ with $s$ uniform outside the
pinned label. Its fifteen nonempty output supports are orthogonal
on the full four-label cube, giving norm $\sqrt{15}r$ before the
same density congruence. Its defect from $J$ is at most $\delta_b$.
A swapped product of two $J+$defect rectangles differs from $J$
by at most $2\delta+\delta^2$. Equality of two uniform coordinates
disintegrates exactly as $q\mathbb E_t$; conditional
Cauchy--Schwarz recovers the original half norms. The sixteen
core pins and four private pins therefore give
$\|D_{\rm sum}-B\|\le E_{\rm pin}$, where $D_{\rm sum}$ is the
sum of those twenty pin terms, with multipliers
$16\cdot4^2$ and $4\cdot16^2$. The condition that the labels involved are distinct is retained
inside every norm estimated above.

The sixteen core/private equalities vanish because $K_C$ is zero
at every core label. Of the other twenty events, compatible
double equalities form matchings on two six-distinct halves.
Every two-coordinate input marginal has point mass at most
$2q^2/(1-q)$: this is attained by a directed core pair; same-side
and private pairs have mass $q^2/(1-q)$, and core/private pairs
have mass $q^2$. The soft density is at most $65536$.
For multiplicity $m$, $m-1\le\binom m2$. Counting even all
$\binom{20}{2}=190$ pairs bounds the positive overlap correction
by $E_{\rm mult}$. Conditioning on the input half contributes the single factor
$h_-^{-1}$ in \eqref{eq:merge-local-I-18}, which also covers every higher coincidence.

With $f_\epsilon(s)=\prod_l(1+\epsilon\chi(b_l-s))$ the
private factor in $F_i^\epsilon$, the exact second moment
$\mathbb E f_\epsilon^2=4(1-2q)$ gives
$\|F_i^\epsilon\|^2=4(1-2q)$, while
$\|G_l\|^2=16\mathbb E\kappa\le16(1+\eta)$.
Hence $\|B\|\le b_B$. In \eqref{eq:merge-local-I-14}, use
$\|w\|,\|v\|,\|e\|\le1$, $\|u\|\le(1-q)^{-1/2}$,
$k\le1$ and $k'\le1+\eta$. The nonconstant protected blocks
are bounded by $A_0,A_1,X$. On the core space orthogonal to
$a$, the constant-label blocks cost at most $16q\omega^2$
and $4q(\eta+\omega)^2$. Thus $P_-BP_-$ (with $P_-$ as defined in the next paragraph) has norm at most
$\alpha$; no auxiliary constant mode is omitted.

Let $S=\operatorname{span}(a\otimes1)$ and $P_-=P\ominus S$.
The old rectangle reduces $S$, $P_-$ and $R$. Its lower diagonal bounds are
$\|a\|_2^2$, $-\|a\|_2\|C_0\|$, and $-\|C_0\|^2$.
After subtracting $B$, the $S$ block is at least $a_-^2-b_B$;
its off-diagonal norm is at most $b_B$. The positive-block elimination in \cref{lem:block-completion}(a)
costs at most $E_S$ times the remaining identity. The other
two blocks have lower bounds $-G,-D$ and off-diagonal norm
at most $b_B$. Since $G>D$, \cref{lem:block-completion}(b) gives their lower bound
$-G-b_B^2/(G-D)$. Adding the two pin errors, compressing to
the physical half and dividing by the old normalizer proves
\eqref{eq:merge-local-I-18}.

For uniform endpoint evaluation, retain certified endpoint
upper bounds on $a_2,c/r,z_2,\varepsilon_Z$. All normalized
numerators are positive increasing in $q,r$, and the lower
denominators are positive decreasing on the gates. Before the
last Schur step replace $G/r,D/r,b_B/r$ by endpoint majorants
$G_0,D_0,B_0$. The lower block matrix with these fixed entries
is valid at every larger prime; use
$G_0+B_0^2/(G_0-D_0)$ after checking $G_0>D_0$.
No monotonicity of a moving gap denominator is assumed.
\end{proof}

\subsection{Allocating the pair-square remainder across a spectral gap}
\label{imp:env:gap-section}

All spaces in this subsection carry probability measures; its $C$ is the rectangle $C_0$ above.  Let $U$ be
the uniform directed $K_{2,2}$ core, with ordered side pairs $a(u),c(u)$,
and let $V$ be the uniform ordered distinct-pair space.  Put
\begin{align*}
 A(u;b)&=K_u(b_0)K_u(b_1),\quad a=A1,\quad C=A-a\otimes1,\\
 a_2&=\|a\|_2,\quad a_\infty=\|a\|_\infty,\quad c=\|C\|.
\end{align*}
Here $K_u$ is the function $K_C$ above with core $C=u$.
The rectangle identities used in
Proposition~\ref{prop:root-chamber} give
$C1=C^*1=C^*a=0$.  Write $M^2=\sup_u\int|C(u;b)|^2\,db$;
the exact row square $256a(u)-a(u)^2$ permits $M^2\le256a_\infty$.
The column square is constant.  Indeed square-affine changes and
private-pair exchange act transitively on $V$ and preserve $C$.
If $F_a,F_c:L^2(V)\to L^2(U)$ are the side-pair pullbacks, the same
symmetry proves
\begin{equation}\label{imp:env:conditional-square}
 F_a^*M_{a^2}F_a=F_c^*M_{a^2}F_c=a_2^2I.
\end{equation}

For real kernels $X,Y:V\to U$, define
\[
 T[X,Y]((u,b),(u',b'))=X(u;b')Y(u';b).
\]
Set
\[
 P_{\rm c}=I_U\otimes J_V,\quad P_{\rm p}=P_a\otimes Q_V,\quad
 P=P_{\rm c}+P_{\rm p},\quad R=Q_a\otimes Q_V,\quad P_a=aa^*/a_2^2.
\]
Put $K=T[A,A]$; expanding $A=a\otimes1+C$ shows that $P,R$ reduce $K$ exactly.
The $R$ block is $T[C,C]$, of norm at most $c^2$.
On $P$, the positive rank-one term $aa^*\otimes J$ and the two adjoint
blocks $a_2C,C^*a_2$ give a lower bound $-a_2c$.

\begin{lemma}[The pair-square residual]
\label{imp:env:pair-residual}
Put $q=p^{-1}$, $r=\sqrt q$, and
$H(b,d)=r^{-1}\mathbb E_s\chi(s-b_0)\chi(s-b_1)
\chi(s-d_0)\chi(s-d_1)$.  Then
\begin{align*}
 \sigma&=(1-2q)^2+\frac{q(1-3q)^2}{1-q},\qquad
 L=H^{\circ2}-\sigma J,\\
 0&\preceq L\preceq\rho I,\quad L1=0,
 \quad\rho=q\left(\frac4{1-q}+2q\right).
\end{align*}
For
\begin{align*}
 E_c&=\tfrac12K\circ\{L(a(u),c(u'))+L(c(u),a(u'))\},\\
 E_b&=K\circ L(b,b'),\qquad E=(2E_c+E_b)/3,
\end{align*}
one has
\[
 PEP\succeq-\alpha P,\quad\|REP\|\le b,\quad
 \|RER\|\le\gamma,
\]
where
\[
 \alpha=\frac{2\rho}{3}\max(a_2^2,c^2),\quad
 \gamma=\rho M^2,\quad
 b=\frac{\rho M}{3}
 \sqrt{(2a_2+c)^2+(a_2+2a_\infty c/a_2)^2}.
\]
The separate $R\leftarrow P_{\rm c}$ bound (the norm of the block $REP_{\rm c}$) is
$b_0=\rho M(2a_2+c)/3$.
\end{lemma}

\begin{proof}
For the raw pair map $W$ with kernel
$\chi(s-b_0)\chi(s-b_1)$, the exact Paley Gram is
$W^*W=qg_eQ+q^2J$, where $g_e=(1-2q-q^2)/(1-q)$.
This gives the constant row sum
$\sigma=(1-q)g_e^2+q^3$.
The auxiliary constant diagonal contributes $(1-2q)^2J$ to
$H^{\circ2}$; the part of $H^{\circ2}$ orthogonal to that contribution is positive with norm at
most $\rho$, by the distinct-pair $K_{2,2}$ bound and the
mean-zero diagonal bound.  Removing its constant eigenvalue proves
the asserted bounds on $L$.

Consider first the unaveraged core placement $L(a(u),c(u'))$.
Its rank/rank term (the part $a\otimes1$ in both factors; the rank/centered and centered/centered terms below are analogous) is supported on $P_{\rm c}$ and costs at most
$\rho a_2^2$, by \eqref{imp:env:conditional-square}.
The rank/centered terms have zero $P$ compression: use $C1=0$
on $P_{\rm c}$, and \eqref{imp:env:conditional-square} followed by
$L1=0$ on $P_{\rm p}$.  Their $R\leftrightarrow P_{\rm c}$ norm is at most
$a_2\rho M$.  To see the normalization explicitly, their squared
Gram is bounded by
$a_2^2\rho^2(CC^*)\circ(F_aF_a^*)$; the last operator is a direct
integral over pair fibers, each with trace at most $M^2$.

The centered/centered term kills $P_{\rm c}$.  Its $P_{\rm p}$ compression is
$B_a^*LB_c/a_2^2$, where $M_a$ is multiplication by the
function $a$, $B_a=F_a^*M_aC$, and $B_c=F_c^*M_aC$.
Equation~\eqref{imp:env:conditional-square} gives
$\|B_a\|,\|B_c\|\le a_2c$, so this compression costs $\rho c^2$.
The full image of $a\otimes g/a_2$ is
\[
 (Cg)(u)(LF_c^*M_aC)(a(u);b)/a_2.
\]
The second factor has constant squared row norm, at most
$\rho^2a_\infty^2M^2$, giving the $R\leftarrow P_{\rm p}$ bound
$\rho a_\infty Mc/a_2$.
Factoring $L$ into two square roots, retaining the same auxiliary
label in the two product maps and the private-coordinate swap,
gives the whole centered/centered norm at most $\rho M^2$.

For the private placement, the rank/rank term is $aa^*\otimes L$
and is positive on $P_{\rm p}$.  The rank/centered term couples $P_{\rm p},R$
with norm at most $a_2\rho M$: its squared Gram is bounded by
$\rho^2(C^*C)\circ I$, whose diagonal is the constant column square.
The centered/centered term kills $P_{\rm p}$, has positive $P_{\rm c}$
compression $CLC^*$, and has $R\leftarrow P_{\rm c}$ norm at most
$\rho Mc$.  The same square-root factorization bounds its full
norm by $\rho M^2$.
Thus the two columns of $REP$ cost respectively
$\rho M(2a_2+c)/3$ and
$\rho M(a_2+2a_\infty c/a_2)/3$.
Their squared sum proves the result.  These are product-operator
bounds, not an arbitrary-kernel Schur-multiplier assertion.
\end{proof}

\begin{proposition}[Finite chamber-cut price for the pair-square law, with the gap bound]
\label{imp:env:finite-gap-chamber}
Use the retained finite constants of
Proposition~\ref{prop:root-chamber}. In its notation (the
following ten aliases are used only in this proposition; in
particular $N,\Gamma,h,\varepsilon,e$ do not have their earlier
meanings here),
$N=N_A$, $\Gamma=\Gamma_{20}$, $E_L=\eps_{\rm lin}$,
$\ell_0=\mathfrak l_0$, $h=h'_p$, $D_{\rm old}=D_{\rm bal}$,
$P_{\rm spike}=\mathfrak s$, $\varepsilon=\eps_Z$,
$h_{\rm mean}=h_0$ and $e=\eps_{\rm sq}$.
Use
$a_-\le a$, $a_+\ge\|a\|_\infty$, $d\ge\|a-1\|_2$,
$\delta=d/\sqrt{1-q}$, and
\[
 \|QL_{\rm soft}Q-rM_{\rm pair}\|\le E_L,\qquad
 \|QL_{\rm soft}1\|\le\ell_0r+E_L,
\]
where $L_{\rm soft}=K\circ\Hbar$ is the soft kernel multiplied by the mean score and
$M_{\rm pair}$ the leading pair matrix of
\cref{prop:root-chamber}.
Let $c$ be an upper bound for $\|C\|$, namely the retained rectangular-domain-gap bound, and assume
$c<a_-$.  Put $d_a=d/a_-$,
\begin{align*}
 E_{0P}&=2rd(\sqrt{30\Gamma}+2)+2a_2N,
 \quad E_{\rm rank}=r(2\delta+\delta^2),\\
 s_+&=1+t^2+2tqr,\qquad s_-=1+t^2(1-4q).
\end{align*}
For a fixed rational $0<\Lambda\le15$ with $k=2t/\Lambda\le s_-$,
let $g(\Lambda)$ be the unique root $>\sqrt{26}$ of
$3g^3+\Lambda g^2-90g-26\Lambda=0$ (the cubic is negative at
$\sqrt{26}$ and strictly increasing beyond it). Define
\begin{align*}
 G&=(s_+-k)a_2c+k\{g(\Lambda)r\Gamma+E_{0P}\}
       +2tE_{\rm rank}+\tfrac{4t}{3}rc^2+t^2\alpha,\\
 d_L&=E_L+rd_a^2/3+qrc^2,\\
 b_L&=E_L+rd_a/3+d_a(\ell_0r+E_L),\\
 h_L&=\ell_0r+E_L(1+d_a),\\
 D&=s_+c^2+2td_L+t^2\gamma+
       \frac{(2th_L+t^2b_0)^2}{s_-a_-^2},
 \qquad B=2tb_L+t^2b.
\end{align*}
Here the residual constants are those of
Lemma~\ref{imp:env:pair-residual}, with $M=16\sqrt{a_+}$ and
$a_2$ in denominators replaced by the smaller bound $a_-$, which only enlarges the constants.
If $G>D$, the independent-half weighted rectangle satisfies
\[
 K\circ\frac13\sum_{j=0}^2(1+t\Hs j)^2
 \succeq-\left(G+\frac{B^2}{G-D}\right)I.
\]
For the actual normalized root, the following number $g$ is
valid in the bound $T\succeq-grI$ for its chamber operator $T$
in the fixed half chart (a \emph{full fixed-half negative
coefficient}):
\begin{equation}\label{imp:env:physical-gap}
 \begin{gathered}
 \frac{h^{-1}[G+B^2/(G-D)]+C_w}{(1-\varepsilon)zr},\\
 C_w=(1+2tB_H+t^2B_H^2)D_{\rm old}
       +\left(\frac{2t}{hr}+\frac{t^2}{hq}\right)P_{\rm spike}.
 \end{gathered}
\end{equation}
Here $B_H=2+r$, $h$, $D_{\rm old}$, $P_{\rm spike}$,
$\varepsilon$, and $z=1+t^2-2t h_{\rm mean}-t^2e$ are the
retained conditioning, collision, spike, normalizer and moment bounds.
The upper bounds for the norm, variance and range of that same
law at the same parameter, proved in
\cref{prop:root-chamber}, remain valid.
\end{proposition}

\begin{proof}
Let $S$ project onto $(a/a_2)\otimes1$, and set $P_-=P-S$.
The exact scalar pair mean is $qr$.  With
$L_0=K\circ\frac13\sum_j(H_j-qr)$,
the weighted rectangle is $sK+2tL_0+t^2E$, where
$s=1+2tqr+t^2\sigma\in[s_-,s_+]$.
Conditional identity \eqref{imp:env:conditional-square} and
the centered pair kernels imply
$SL_0S=SES=SL_0P_-=SEP_-=0$.
The positive $S$ block is therefore $sa_2^2$.

On $P_-$, the exact old bound is $-a_2c$.
In comparing $K$ with its leading stars, $T[C,C]$ has zero
$P$ compression, so only the rank/centered remainder $E_{0P}$
is charged.  The inherited eighteen-feature Gram bound $\Gamma$
also applies on the independent half: cross-factor Grams vanish,
and within-core and within-private conditional counts dominate
the smaller counts used in its original proof.
The finite leading matrix has negative norm $g(\Lambda)$;
its remaining blocks cost at most
$\max(2,\Lambda/3)\le5<g(\Lambda)$.
The exact positive-coefficient splitting
\[
 sK+2trM_{\rm pair}=(s-k)K+k(K+\Lambda rM_{\rm pair})
\]
therefore incurs the error $E_{0P}$ only in the second summand.
The weighted rank error is $E_{\rm rank}$.
The other core $P_{\rm p}$ linear block costs at most $2rc^2/3$;
the private $P_{\rm c}$ block is positive.  This proves the stated $G$.

The angle of $P_-$ to the ordinary constant is at most $d_a$.
Splitting through $Q$ gives $b_L$; the private leading pair term
has coefficient $1/3$, while core terms have private-constant output
and vanish against $R$.  The same split gives $d_L$, with $qrc^2$
paying the removed raw mean.  For the $S$ column, the leading pair
matrix sends $Q((a/a_2)\otimes1)$ into $U\otimes1$, which $R$
annihilates; this gives the sharper $h_L$.
Eliminating the positive $S$ block costs only
$(2th_L+t^2b_0)^2/(s_-a_-^2)$ on $R$.
The remaining two-block bound is $G+B^2/(G-D)$ by
\cref{lem:block-completion}(b).
It is already a full-space bound, so no second constant-block
penalty is added after compression.

Actual-half compression costs $h^{-1}$.
Outside the configurations where the two label pairs of a score
are equal as unordered sets, the weight is at most
$1+2tB_H+t^2B_H^2$; the equal-pair strata are paid by the
two explicit spike terms in \eqref{imp:env:physical-gap}.
The retained raw mass is at least $(1-\varepsilon)z$.
These facts prove the bound in the original twelve labels, in
the fixed half chart, without passing to the tilted half marginal.

For a uniform endpoint price, majorize $G/r,D/r,B/r$ first by
fixed endpoint values $G_0,D_0,B_0$ and check $G_0>D_0$.
Use $G_0+B_0^2/(G_0-D_0)$ at every larger prime.
All constituent normalized numerators increase in $q,r$, and
all positive lower denominators decrease.  This order of operations
avoids assuming monotonicity of a quotient with a moving gap.
The fixed $\Lambda$ remains feasible: at the endpoint $X$ (the
lower endpoint of the interval) $k\le s_-(X)$, and $s_-(p)$ is
nondecreasing in $p$, so $k\le s_-(p)$ for all $p\ge X$.
\end{proof}

\begin{corollary}[Interpolation after physical normalization]
\label{imp:env:physical-interpolation}
Let $g_0r$ be any full negative price for the mass-one unchanged
physical root $\Gamma_0$.  Fix $0<\tau\le1$.  In
Proposition~\ref{imp:env:finite-gap-chamber}, replace
\[
 s_+=\tau+t^2+2tqr,\qquad s_-=\tau+t^2(1-4q),
\]
and replace the constant $1$ in the collision coefficient
$1+2tB_H+t^2B_H^2$ by $\tau$.  Keep all linear and square
coefficients $2t,t^2$, and check the same positive-block and gap
conditions. After these replacements, choose endpoint majorants
$G_0,D_0,B_0$ for $G/r,D/r,B/r$, respectively, with $G_0>D_0$, as
in the final paragraph of the proposition's proof, and set
\[
 C_{\rm soft}=G_0+\frac{B_0^2}{G_0-D_0}.
\]
This is the raw soft coefficient, before half-conditioning, scalar
normalization and collision corrections. If $C_{w,\tau}$ is the
corresponding collision price in \eqref{imp:env:physical-gap}, then
\[
 g_{\rm int}=
 \frac{(1-\tau)g_0+
       [C_{\rm soft}/h+C_{w,\tau}/r]/(1-\varepsilon)}{z}
\]
is a full negative coefficient, in the sense above
($T\succeq-g_{\rm int}rI$), for the original individual-square
(pair-score) law. Here $h$, $\varepsilon$ and $z$ are those of the proposition at the same $t$.
Its marginal and layer data remain those at the original parameter $t$.
\end{corollary}

\begin{proof}
Before normalizing the new tilt, split its exact physical operator as
\[
 (1-\tau)\Gamma_0+
 \Gamma_0\left(\tau+2t\,\frac13\sum_jH_j
                       +t^2\,\frac13\sum_jH_j^2\right).
\]
The first summand has lower price $(1-\tau)g_0r$ in the mass-one
physical convention.  The second summand can have signed entries;
the three-block argument requires only its positive $S$ block and
the stated gap gates.  Its raw scalar normalization alone costs
$(1-\varepsilon)^{-1}$.  Absolute collision weights are bounded by
$\tau+2tB_H+t^2B_H^2$, with the same exceptional-pair terms.
Add these two lower prices and divide once by the common new mass
lower bound $z$.  No second old-normalizer or half-conditioning
factor belongs to the first summand.
\end{proof}

\subsection{Retaining the constant coefficient at the layer cut}
\label{imp:env:layer-section}

\begin{proposition}[Layer-cut remainder, summed without expanding all monomials]
\label{imp:env:layer-resummation}
For the individual-square law in Proposition~\ref{prop:root-layer}, put
\begin{align*}
 d_\mu&=\frac8{1-q^2},\quad M_\Phi=d_\mu(1+q)^3,
 \quad b_*=6+51d_\mu,\\
 n&=r\{\kappa(1+5r)+b_*r\},\quad
 \kappa=(\sqrt3+\sqrt7)/2,\quad B_H=2+r.
\end{align*}
Replace its remainder coefficients by
\begin{align*}
 C_{\rm nn}&=(n/r)^2+3969\cdot16(2tB_H+t^2B_H^2),\\
 C_0&=128M_\Phi+8d_\mu^2(1+q)^5,\\
 C_{\rm sc}&=C_0+127(20tr+16t^2q+2t^2).
\end{align*}
Keep $C_{\rm cr},C_{\rm in}$ and both mass lower bounds unchanged,
and substitute $C_{\rm nn}+C_{\rm sc}+C_{\rm cr}$ in its negative,
centered and marginal formulas.  The resulting prices hold for
the same positive law, with no change to its support or parameter.
\end{proposition}

\begin{proof}
Write the probability-normalized triangle operator of \cref{prop:root-layer} as $\Phi=J+N$,
where $\|N\|\le n$.  In the score-independent coefficient retain
the full $N\otimes N$, of norm at most $n^2$, instead of expanding
its $63^2$ nonempty monomials.  The score-containing terms retain
their specified four-label mask, with all labels distinct, and the two-Paley
price, giving $C_{\rm nn}$.

Delete the removed constant mask from the whole positive
$\Phi_1\Phi_2$ of the two triangle operators.  Among its four specified equalities, two twins
cost $64M_\Phi q$ each and two cross-component pins cost
$4d_\mu^2(1+q)^5q$ each.  Their positive union bound gives $C_0$,
including all intersections.  For terms containing a score, single
coincidences in the specified single strata have
$|H|\le2r$ for the pair score $H$ of
\cref{def:positive-root}, and equal unordered pairs
have $|H|\le1/r$; the unchanged mass bounds $4q,2q^2$ and
absolute base density $127$ give the remaining terms of $C_{\rm sc}$.
In particular no mask is deleted through a signed $N\otimes N$
kernel.  Its row contribution is bounded by its full norm, so the
same replacement is valid in the marginal estimate.
The input-half conditioning factor from \cref{prop:root-layer} remains an upper bound here.

The saving in the remainder coefficient is exactly
\[
 3969\cdot16-(n/r)^2+127(4+2q)-C_0,
\]
whose limit is $62476-\kappa^2>62471$.
After cancellation of $r$, every displayed endpoint expression has
positive increasing numerators and decreasing positive lower
denominators.  The explicit mass, centered-gap and tensor-domain
gates (\cref{app:certificates}) therefore certify uniform prices above the endpoint.
They must still be combined with a chamber price for the same $t$,
fresh native completion, and the corresponding marginal restoration.
\end{proof}


\subsection{Supports of the $39$-trace law}

The tilted law of the last interval has weight
$\eta+(1+tH)^2$, where each trace in $H$ uses the same auxiliary
summation label on its two halves. Squaring introduces two auxiliary
labels; their diagonal must be retained. The following lemma
bounds the product synthesis on one half, including this diagonal,
before it is contracted with the other half.
\begin{lemma}[Higher-star products]
\label{imp:local:star-product}
Let $H=L^2_0(\mathbb F_p)$ and
$U_S=(\psi_S-\mathbb E\psi_S)/r:H\to L^2(C)$, where
$\psi_S(s)=\prod_{i\in S}\chi(s-C_i)$.
For subsets $S,T$ of the four core coordinates with $S$ nonempty and
$|T|\ge2$, let $U_S\cdot U_T:H\otimes H\to L^2(C)$ be the
pointwise product synthesis, defined on simple tensors by
$(U_S\cdot U_T)(f\otimes g)=(U_Sf)(U_Tg)$.
It has the decomposition
\begin{equation}\label{eq:merge-local-I-4}
 \begin{aligned}
 U_S\cdot U_T&=h e_0^*+R,\quad Re_0=0,\quad \|R\|=O(1),\\
 h&=1_{S=T}/r+O(1)\ \text{pointwise},\qquad
 e_0(s,u)=\frac{p1_{s=u}-1}{\sqrt{p-1}}.
 \end{aligned}
\end{equation}
Here $h$ is a function on core configurations and the $O(1)$
constants are absolute, uniformly in $p$ and in the labels. The
same conclusion holds for fixed linear combinations, with the
coefficient of $1/r$ equal to their coefficient-vector inner product.
A coordinate evaluation times a higher star has product norm bounded by an absolute constant.
The assertion does not include a shared singleton--singleton support.
\end{lemma}

\begin{proof}
For $Df(s,u)=\sqrt p1_{s=u}f(s)$, its projection $D_H$ to
$H\otimes H$ obeys
\begin{equation}\label{eq:merge-local-I-5}
 D_H^*D_H=(1-2q)P_H+(1-q)J.
\end{equation}
On the orthogonal complement of its range, auxiliary functions vanish
on $s=u$ and all auxiliary centering terms disappear.  If $S\ne T$,
choose two output coordinates whose incidence with $s,u$ is a
matching or a three-edge path.  Such coordinates exist also when one
support contains the other.  Two Paley transforms give norm $q$,
with a bounded middle multiplication for the path.  If $S=T$, choose
two common coordinates and use the restricted $K_{2,2}$ estimate
$2q$, retaining distinctness of both output and auxiliary pairs.
After fixing all other core coordinates, their characters are bounded
auxiliary multipliers.  Exclusions of fixed coordinates are separate
output restrictions.  The actual core density is
$16/\theta$, $\theta=(1-q)(1-3q)(1-7q)$, so output compression costs
at most $4/\sqrt\theta$.  Division by $r^2=q$ proves the off-diagonal
bound without deleting a coupled injectivity mask.

On the mean-zero part of $D_H$, repeated-root cancellation gives
$U_{S\triangle T}$, or zero if $S\triangle T$ is empty, plus at most
four evaluation terms of size $q/r=r$.  Terms containing an auxiliary
mean are bounded by the degree-four character estimate (\cref{thm:weil}).  Equation
\eqref{eq:merge-local-I-5} gives a bounded diagonal inverse.  Finally the $e_0$ image is
\[
 \frac{\mathbb E_s[(\psi_S-m_S)(\psi_T-m_T)]}
      {r\sqrt{1-q}},\qquad m_S=\mathbb E_s\psi_S.
\]
Its numerator is $1+O(q)$ when $S=T$ and $O(r)$ otherwise.
For evaluation times a higher star, fix the evaluated coordinate;
one other leaf supplies a Paley factor $r$, canceling the star's
normalization. Passing from the product law to the actual law
changes the bound only by the density factor $\prod_jD_j$ of
\cref{lem:l20-star-matrix}, an absolute constant.
\end{proof}

\subsection{The $39$-trace law at every sorting cut}
\label{imp:local:filtered-coherent}

We now define the list of supports. A symmetrised list
$\mathcal D_{51}$ of $51$ supports is defined first, and the
action of its score and of its square on the soft kernel (the
kernel without the distinctness indicator) is computed, keeping
every mean term of a star, including the functions invariant under
simultaneous translation of all labels. Removing the twelve
supports with fewer than three labels in one layer then leaves
the list $\mathcal D_{39}$ used in the law, and explicit finite
matrix bounds at the endpoint parameter $t_*=12/5$ follow.

The twelve vertices are grouped into six bases, numbered
$0,\ldots,5$, each carrying two labels $(j,\ell)$,
$\ell\in\{0,1\}$, in the two layers (the ``twins'' $(j,0),(j,1)$).
The directed edges between the bases are
\begin{equation}\label{eq:merge-local-I-19}
 0\longrightarrow2\longrightarrow4\longrightarrow0,\qquad
 1\longrightarrow5\longrightarrow3\longrightarrow1.
\end{equation}
At the first chamber cut, with the two halves identified by the
reflection (the reflected chart), the halves are
$(a,b,c)=(0,1,2)$ and $(a',b',c')=(5,4,3)$.
On each four-label core retain the four cubic supports and the
quartic support. Give a support $S$ sign $(-1)^{|S\cap a|}$.
Form the twenty-five products of pairs of supports, one from
each core, divide their coefficients by five, and average over
\begin{equation}\label{eq:merge-local-I-20}
 j\mapsto j,\ j+2,\ j+4,\ 1-j,\ 3-j,\ 5-j\pmod6.
\end{equation}
After collecting equal physical supports this gives seventy-five
nonzero coefficients, each $\pm1/15$.

Retain a support if, at every chamber cut, either its
intersections with the core labels of both halves have size at
least two, or its intersections with the private labels of both
halves have size two. Denote this fifty-one-element
dictionary by $\mathcal D_{51}$, retaining the collected
coefficients $c_D$. Put
\begin{equation}\label{eq:merge-local-I-21}
 H_{\mathcal D}
   =\frac1r\sum_{D\in\mathcal D}c_D\,
                \mathbb E_s\prod_{i\in D}\chi(s-x_i),
 \qquad \sigma_{\mathcal D}=\sum_D c_D^2 .
\end{equation}
Here $x=(x_i)$ is the configuration of the twelve labels.
This definition uses physical labels and is independent of a
chosen cut. In particular $\sigma_{51}=17/75$ and
$\sum|c_D|=17/5$.
We give the complete argument keeping the common auxiliary
label $s$ (the root) unintegrated, since integrating it out
first and using only moment orthogonality would not suffice. Write
$K_{\rm soft}=T[A,A]$, with $A$ the exact swapped rectangle of \cref{prop:l20-rectangular}, and
\begin{equation}\label{eq:merge-local-I-23}
 A=a\otimes1+r(\sqrt{30}\,Zv^*+2we^*)+E,\qquad
 \|E\|=O(q),\quad a=1+O(r).
\end{equation}
Both row and column squared norms of $A$ are at most
$256a_\infty$ (with $a_\infty=\|a\|_\infty$). For two trace supports, retain their common
auxiliary labels $s,t$ and split $s=t$ from $s\ne t$.
If two nonempty raw supports have union of size at least two,
their off-diagonal two-auxiliary map has norm $O(q)$.
Distinct supports use a matching or three-edge path; equal
supports use the restricted $K_{2,2}$ estimate. Output and
auxiliary distinctness are both retained.

If the two maps obtained from the core part of $A$ have the
property of the preceding paragraph (norm $O(r)$ after one free
transform) on both halves, write the modulated rectangles
(rectangles $A$ with left and right multiplication operators
$M_f,M_g$ by the displayed private multiplier functions $f,g$
and their reflected counterparts $\tilde f,\tilde g$) as
$B_{st}=M_fAM_g$ and
$D_{st}=M_{\tilde f}AM_{\tilde g}$. The private multipliers have
modulus at most one, so positive Gram Schur order gives
\begin{equation}\label{eq:merge-local-I-24}
 \left\|\mathbb E_{s\ne t}B_{st}B_{st}^*\right\|
 \le256a_\infty\|F_{\rm core}\|^2=O(q^2),
\end{equation}
where $F_{\rm core}$ is the core synthesis operator of
\eqref{eq:merge-local-I-23}; and the same for $D$. Operator Cauchy--Schwarz bounds
$\mathbb E B_{st}\otimes D_{st}^*$ by $O(q^2)$, with the private
swap unchanged. Division by the score-square normalization $q$
gives $O(q)$. If both private maps instead have the property,
use the input Grams and the column bound. Exact support
enumeration covers all $51^2$ ordered pairs: 2537 use core output
Grams and 64 use private input Grams.

On $s=t$, cancel repeated roots. The identical-support term is
$K_{\rm soft}+O(q)$. Other terms are raw traces on the symmetric
difference, with at most twelve hole terms (terms where the
common auxiliary label equals a fixed label) of size $q$.
A raw trace with nonempty supports on both halves costs $O(q)$
after multiplication by $K_{\rm soft}$. For core/core or
private/private placement, use the one-auxiliary form of
\eqref{eq:merge-local-I-24}, with two raw norms $O(r)$. Pure core/private placement
reduces to $A$ times a raw-star cross-kernel; expand \eqref{eq:merge-local-I-23}.
Rank and remainder terms cost $O(q)$. Every leading edge has
a coordinate evaluation at one endpoint. The higher/raw
product has a bounded exceptional $e_0$ image and an $O(r)$
remainder, whereas the evaluation/raw image is
\begin{equation}\label{eq:merge-local-I-25}
 \frac{\psi_S(x_i)-\mathbb E\psi_S}{\sqrt{p-1}}=O(r).
\end{equation}
The old factor $r$ makes these terms $O(q)$. A shared singleton
pair gives $1_{x_i=y_j}-q$, whose positive equality mass is
$O(q)$. One-sided differences give invariant functions of size
$O(r)$; their leading one-sided ranks vanish under ordinary
centering, and multiplication by $K_{\rm soft}-J=O(r)$ costs
$O(q)$. Thus
\begin{equation}\label{eq:merge-local-I-26}
 Q(K_{\rm soft}\odot H^2)Q=\sigma QK_{\rm soft}Q+O(q).
\end{equation}

The corresponding linear calculation gives
\begin{equation}\label{eq:merge-local-I-27}
 Q(K_{\rm soft}\odot H)Q=QHQ+O(q),\qquad
 \|K_{\rm soft}\odot H\|=O(r).
\end{equation}
Here are the exceptional shared-label cases. Rank/star terms
have a free raw core and private transform, each $O(r)$.
For the product of two stars, the two arms (synthesis factors)
from the untilted law are $(Z,v)$ and $(w,e)$.
The cross pairing has a free leaf at its $(w,v)$ endpoint,
since each half score uses at least three labels.
For equal $(Z,v)$ pairings, two private leaves again give a free
transform. Otherwise the private support is a singleton and
the core support has at least two labels. Split its higher/raw
product into the constant exceptional image plus an $O(r)$
remainder. The two exceptional images contract against the
swapped private evaluations; the delta terms vanish by
$\chi(0)=0$ and the remaining probability operator is
$qQ/(p-1)$. It is $O(q^2)$, not an omitted order-$r$ block.
The equal $(w,e)$ case exchanges core and private. Every
single remainder costs $O(q)$. An $E$/star term applies $E$
fiberwise between two bounded product maps with the shared
label retained. The $E$/rank term has free raw transforms;
$E$/$E$ costs $O(\|E\|^2/r)$. These prove \eqref{eq:merge-local-I-27}.

For a half support $S$ (a support contained in one half), extending the notation of \cref{imp:local:star-product}, let
$U_S=(\psi_S-m_S)/r$ on
mean-zero auxiliary labels and $M_S=m_S/r$. With $C$ the
symmetric matrix, indexed by the half supports, whose entries
are the collected coefficients $c_D$ of
\eqref{eq:merge-local-I-21}, and $U$ the row of operators
$(U_S)_S$,
\begin{equation}\label{eq:merge-local-I-28}
 H=rUCU^*+rMCM^* .
\end{equation}
The two terms occupy respectively nontrivial translation sectors
and the invariant sector. To specify the finite matrix gate
without ambiguity, index all nonempty half supports by $S$.
Let $D_{ij}$ be $+1$ or $-1$ according to the directed cross
edge between coordinate $i$ and the reflected coordinate $j$,
and zero if it is absent. The old matrix has singleton block
$iD$, higher/singleton entry $\prod_{i\in S}D_{ij}$, its
adjoint entry, and zero higher/higher block.

The Gram matrix of the core raw mean functions $M_S$ is at most $16+O(r)$, by cube orthogonality
and the core density. Group mixed means by their private singleton
or pair support. Their raw Gram is at most
$q[1+q/(1-q)+2qr/(1-q)]I$; for each fixed injective core the
character supports have Gram $I+O(r)$. Conditional synthesis gives
mixed mean Gram $I+O(r)$. The Gram matrix of the common actual stars is likewise
$I+O(r)$ on nonconstant auxiliary labels. These estimates will price
the invariant sector as well as the nonzero translation sectors.

\paragraph{Physical completion.}
Each half support has at least three labels. Complete character
cancellation therefore requires at least three cross identifications, of positive row/column mass $O(q^3)$.
Everywhere each single-support trace factor in \eqref{eq:merge-local-I-21}, before its coefficient, is at most $1/r$ in absolute value, so even at the configurations where
a trace takes an exceptional value (the spikes), its square
contributes $O(q^2)$ after the normalizations. Other collisions have mass $O(q)$ and bounded
score by the repeated-root Weil bound (\cref{thm:weil}). This controls all
physical deletions and their intersections.

The positive constant $\eta=1/100$ added in the weight (its
floor) and the Weil bound on the strict support (distinct
labels) give uniformly
positive bounded new marginal densities. The full identities for the new half-marginal density $g$
with respect to the fixed half law give $\|g-1\|_2=O(r)$. For each normalized star, the weighted
affine Gram with weight $(g^{-1/2}-1)^2$ is real symmetric
and a scalar multiple of the identity on mean-zero labels,
with scalar its trace divided by $p-1$; this shows that multiplying the finite feature synthesis by
$g^{-1/2}-1$ changes its norm by $O(r)$. Thus the leading
$r$ operator changes by $O(q)$ under actual-marginal
congruence. The finitely many invariant mean arms are pointwise
bounded; their multiplication by $g^{-1/2}-1$ likewise has norm
$O(r)$ by its $L^2$ bound, covering the second term of \eqref{eq:merge-local-I-28}.
Finally, writing $T$ for the chamber operator of the tilted law
on the induced half,
\begin{equation}\label{eq:merge-local-I-29}
 T=g\otimes g+QTQ-(g-1)\otimes(g-1)
\end{equation}
keeps the constant rank positive and charges only $O(q)$.
This proves that the leading bound transfers from the fixed half chart to the actual half marginal at every chamber cut.

Remove from $\mathcal D_{51}$ the twelve supports containing
fewer than three labels in one layer. Call the remaining
dictionary $\mathcal D_{39}$. It has twelve supports of degree
six, twenty-four of degree seven, and three of degree eight.
Consequently
\begin{equation}\label{eq:merge-local-I-30}
 \sigma_{39}=13/75,\quad \sum|c_D|=13/5,\quad
 |H_{\mathcal D_{39}}|\le14+r
 \quad\hbox{on strict physical support}.
\end{equation}
The last bound uses $(k-2)r+q$ for even degree $k$ and
$(k-1)r$ for odd degree, before dividing by $r$.

\begin{lemma}[Admissibility of the coherent root]
\label{lem:imp:coherent39-law}
For every rational $t\ge0$ and every $\eta>0$, normalize
$\eta+(1+tH_{\mathcal D_{39}})^2$ against the old physical root at the prime $p\equiv3\pmod4$ fixed above.
The resulting probability law has exactly the old root's support
and accepts the six queries. It is invariant under the three
chamber reflections $j\mapsto2a-1-j$, $a=0,1,2$, the rotation
$j\mapsto j+2$, the simultaneous layer interchange, and the
square-affine maps of the labels. Every slot is uniform.
These assertions also hold for the primewise product on a block.
In particular they hold at the selected parameter $t=67/40$ and
floor $\eta=1/100$.
\end{lemma}

\begin{proof}
The initial average in \eqref{eq:merge-local-I-20} makes the signed
support dictionary invariant under the chamber reflections and
the rotation. Simultaneous layer interchange permutes the four
cubic supports and fixes the quartic support in each core; it
preserves their signs and the twenty-five corresponding products.
The filter defining $\mathcal D_{51}$ is imposed at every chamber
and is unchanged by these permutations. The subsequent filter
requiring at least three labels in each layer is also unchanged.
Thus each of these permutations preserves the collected
coefficients of $\mathcal D_{39}$ and fixes its score.

Under $x_i\mapsto\alpha x_i+\beta$, with $\chi(\alpha)=1$,
the substitution $s=\alpha s'+\beta$ fixes every character
product in \eqref{eq:merge-local-I-21}, in both odd and even
degrees. All the stated maps preserve the old physical root,
so they preserve its tilted weight and normalizer. Since the
weight is at least $\eta>0$, normalization is defined and the
support is unchanged; query acceptance follows from that of
the old root. Translations act transitively on the labels,
so every slot is uniform. Tensor products give the block
assertions. Only the simultaneous layer interchange is needed;
independent flips of individual twins are not asserted.
\end{proof}

\begin{proposition}[Thirty-nine coherent traces at every sorting cut: the leading coefficients]
\label{imp:local:coherent-39}
Let $\Gamma$ be the coherent law at the endpoint parameter
$t_*=12/5$, that is, the uniform law reweighted by
$1/100+(1+(12/5)H_{\mathcal D_{39}})^2$ and normalized. It is
invariant under the chamber reflections
$j\mapsto2a-1-j$, $a=0,1,2$, and under the simultaneous
interchange of the two layers $(j,\ell)\mapsto(j,1-\ell)$; these
are the reflections occurring in the sorting sequences.
(Independent flips of single twins are not symmetries and are not
used.) There are absolute constants $C$ and $P_0$ such that for every
prime $p\equiv3\pmod4$ with $p\ge P_0$ the Markov operators
of $\Gamma$ at its cuts, on its actual marginals, obey
\begin{equation}\label{eq:merge-local-I-31}
 T_{\rm chamber}\succeq-(5.413r+Cq)I,\qquad
 T_{\rm layer}\succeq-(2.19r+Cq)I,
\end{equation}
the first at all three chamber cuts. This proposition identifies
only the leading coefficients $5.413$ and $2.19$; by itself it
gives no bound at any specified prime. The explicit bounds used
in the construction, with every remainder evaluated, are
\cref{imp:local:finite-coherent-layer,imp:env:finite-coherent-chamber}.
\end{proposition}

\begin{proof}
The chamber calculation for $\mathcal D_{51}$ above applies to the retained dictionary $\mathcal D_{39}$.
There are 1493 core-output and 28 private-input covariance
cases. Its strength is
$\lambda=2t/(1+\eta+t^2\sigma)=12000/5021$.
Exact rational gates (\cref{app:certificates}) give chamber coefficient $5413/1000$,
core mean coefficient $7/100$ and mixed mean coefficient
$21/100$. The invariant chamber price is below
$\lambda(16\cdot7/100+21/100)<16/5$.

At the layer cut a half is initially two independent directed
cyclic triangles, on bases $\mathcal A=\{0,2,4\}$ and
$\mathcal B=\{1,3,5\}$. The actual half conditions their six
labels to be distinct, at cost $1+O(q)$. Each component has
\begin{equation}\label{eq:merge-local-I-32}
 \Phi=J+rFM_0F^*+E,\quad \|E\|=O(q),\qquad
 M_0=\begin{pmatrix}iD_0&-I_3\\-I_3&0\end{pmatrix},
\end{equation}
where $F$ is the synthesis operator assembled from the six
feature maps of the preceding paragraph. Its arms (factors) are
the three evaluations and three pair stars.
The six omitted pair means give $+6qJ$, included in $E$.
There is no higher/higher edge.

For every retained trace each half support meets both
triangles in one or two coordinates and has size at least
three. In the quadratic product use
$q^{-1}\mathbb E_{s\ne t}B_{st}\otimes D_{st}$.
Operator Cauchy--Schwarz now uses $BB^*$ and $D^*D$:
the small raw maps must be output $\mathcal A$ and input
$\mathcal B$, or input $\mathcal A$ and output $\mathcal B$.
The exact enumeration gives 1437 cases of the first kind and 84
of the second. Each raw map is $O(q)$; the component density
$8/(1-q^2)$ and bounded row/column squares of $\Phi$ give
the full $O(q)$ off-diagonal contribution.

On the auxiliary diagonal, supports in opposite components
use the one-label crossed argument with two $O(r)$ maps.
If both lie in the same component, expand \eqref{eq:merge-local-I-32}.
The rank and remainder terms cost $O(q)$; each star has an
evaluation endpoint, so \eqref{eq:merge-local-I-25} controls the exceptional
image. Two singleton supports instead give an equality
kernel of bounded-density mass $O(q)$.
For this particular dictionary, collecting all one-sided
symmetric differences gives thirty pair supports of total
absolute coefficient mass $8/25$. After conditioning the
half to be six-distinct, every such raw trace is exactly
$-q$. Therefore at either cut the full covariance, in the
fixed induced-half chart, is, writing $K$ for the tensor square
of the layer-cut kernel,
\begin{equation}\label{eq:merge-local-I-33}
 K\odot H^2=\sigma K+O(q).
\end{equation}
This full-norm identity is not asserted on the independent
triangle chart before conditioning.

For the layer linear term expand the two copies of \eqref{eq:merge-local-I-32}.
One star term and one rank term have, on both sides, free raw
star factors carrying none of the displayed characters (passive
factors). For two stars, the scalar is $r^2/r=r$.
A coordinate/coordinate endpoint has a free leaf because
the score uses at least three coordinates. Every
higher/higher endpoint is opposite one of this type.
The remaining coordinate/higher against higher/coordinate
case either has a free leaf or has two identical pair-star
exceptional images. The latter is the same shared-label
$qQ/(p-1)$ contraction used above; all single remainders
have norm $O(r)$. An $E$/star term is a fiberwise $E$
between bounded product maps; $E$/rank has free raw
transforms and $E$/$E$ costs $O(\|E\|^2/r)$.
Hence $K\odot H=H+O(q)$ on the induced half.

The real vector space spanned by the real and imaginary parts
of the layer feature functions that occur (the active layer
realification) has dimension $54$; rational
elimination proves coefficient $219/100$. Its raw-mean
coefficient matrix is bounded below by $-21/100$.
Group mean arms by the singleton or pair in one triangle.
The normalized six-arm nonconstant Gram is $I+O(r)$;
the three pair constant images add only $3q^2J$ before
normalization. Fixed opposite-triangle character Grams
are $I+O(r)$, so the complete mean frame is $I+O(r)$.
Its invariant negative coefficient is below
$\lambda(21/100)<1$. Thus no invariant sector is omitted.
Physical collision deletion, positivity and actual-marginal
transport follow exactly as in the physical-completion argument above:
all retained half supports still have size at least three.
The dictionary, all 1521 crossed cases, chamber and layer
matrix gates and separate mean gates are reconstructed by
\texttt{check\_coherent\_39\_all\_cuts.py} (\cref{app:certificates}).
\end{proof}

\subsection{One parameter $\ulaw$ for both cuts of the $39$-trace law}

The matrix bounds above were proved at the endpoint parameter
$t_*=12/5$, because the finite checks are cleanest there. The law
used in the construction has the smaller parameter $t=67/40$. The
following proposition transfers the endpoint bounds to every
$0\le t\le t_*$ by convexity---the leading matrix is affine in an
explicit interpolation parameter $\alpha(t)$---at both the
chamber cut and the layer cut simultaneously; all remaining terms
are then evaluated directly at $t=67/40$.

\begin{proposition}[The common coherent parameter interval]
\label{imp:local:coherent-parameter}
\label{imp:env:coherent-parameter}
Keep $\eta=1/100$, $\sigma=13/75$ and the 39-trace score, and let
$0\le t\le t_*=12/5$. Set
\begin{equation}\label{eq:merge-local-I-43}
 z(t)=1+\eta+\sigma t^2,\qquad
 \lambda(t)=\frac{2t}{z(t)},\qquad
 \alpha(t)=\frac{\lambda(t)}{12000/5021}.
\end{equation}
Then $0\le\alpha(t)\le1$. In the leading matrices of the normalized
coherent law, valid negative coefficients are
\[
 g_{\rm ch}(t)=(1-\alpha(t))\sqrt{30}
                   +\alpha(t)\frac{5413}{1000},\qquad
 g_{\rm lay}(t)=\frac{219}{100}.
\]
The invariant mean branches obey the same displayed bounds.
Each estimate in \cref{imp:local:finite-coherent-layer,imp:env:finite-coherent-chamber}
below applies to the physical weight
$\eta+(1+tH)^2$ by using this $t$, its scalar leading mass $z(t)$,
and the corresponding leading coefficient. The error terms in the linear and covariance identities for
the score itself do not change with $t$. A separately certified matrix bound
at a specified parameter may improve its leading coefficient, provided
the separate invariant-mean price remains below it. This changes neither
the law nor its other prices.
\end{proposition}

\begin{proof}
The exact factorization
\begin{equation}\label{eq:merge-local-I-44}
 t_*z(t)-tz(t_*)=(t_*-t)(1+\eta-\sigma tt_*),\qquad
 1+\eta-\sigma t_*^2=29/2500
\end{equation}
gives $0\le\alpha\le1$. At either cut, with $M_{\rm old}$ and $C$ denoting the untilted leading matrix and the score coefficient matrix at that cut, the leading coefficient matrix is
\[
 M_{\rm old}+\lambda(t)C
 =(1-\alpha)M_{\rm old}
       +\alpha\bigl(M_{\rm old}+(12000/5021)C\bigr).
\]
The endpoint orders are those proved in
\cref{imp:local:coherent-39}. The old chamber coefficient is $\sqrt{30}$;
the old layer coefficient $(\sqrt3+\sqrt7)/2$ is less than $219/100$.
Convexity gives the stated bounds. The old invariant centered leading
operator is zero, so the separate invariant branches interpolate as well.

For the finite calculation, expand the raw weight before normalizing.
The old-root, linear and covariance terms have coefficients
$z(t),2t,t^2$, respectively; the regular and exceptional collision
weights are their explicit positive quadratic polynomials. Each
primitive error is bounded separately below. Substituting these
coefficients therefore prices the same physical law at both cuts.
The $10q$ allowance for centering the layer mean arms is bounded at
$t_*$ and remains valid for smaller $t$ (see the proof of \cref{imp:local:finite-coherent-layer}). For each fixed $t$, the
positive endpoint majorants give uniformity in $p$; no monotonicity
of the final price in $t$ is needed.
\end{proof}

\subsection{An explicit layer price for the coherent law}
\label{imp:local:finite-coherent-layer-section}

We now evaluate the layer-cut bound for the same $39$-trace law
numerically. All norms in this subsection are taken in $L^2$ of
the fixed uniform half law (the fixed chart), not in the tilted
law's own marginal; the passage from the chart to the actual
marginal was already bounded in the transport estimates above, so
no extra factor for that change of measure is charged here.
Fix $0\le t\le12/5$ and keep $\eta=1/100$, $L=13/5$,
$\sigma=13/75$ and $z=z(t)=1+\eta+t^2\sigma$ from
\cref{imp:local:coherent-parameter}.
For $p\ge2^{32}$ set
\begin{equation}\label{eq:merge-local-I-34}
 \begin{aligned}
 h&=1-9q,& d&=\frac8{1-q^2},&
 M&=d(1+q)^3,& \Gamma&=1+1000r,\\
 B_H&=14+r,& W&=\eta+(1+tB_H)^2 .
 \end{aligned}
\end{equation}
Define positive coefficients
\begin{equation}\label{eq:merge-local-I-35}
 \begin{aligned}
 C_{\rm soft}
   &=\frac{1100z+2tL\cdot100000+t^2\cdot200000+10}{h},\\
 C_\times&=\frac W h[384M+72d^2(1+q)^5],\\
 C_{\rm exc}
   &=\frac{9830400d^2}{h}
          [(1+\eta)q^2+2tLqr+t^2L^2q],\qquad
 C=C_{\rm soft}+C_\times+C_{\rm exc}.
 \end{aligned}
\end{equation}
Here $C$ is a finite coefficient, not an unspecified asymptotic
constant. Put
\begin{equation}\label{eq:merge-local-I-36}
 m=z-Cq,\qquad
 \delta=\frac{Cq}{m},\qquad
 c=\frac{4z\Gamma r+Cq}{m}.
\end{equation}

\begin{proposition}[Finite layer-cut bound for the coherent law in the fixed half chart]
\label{imp:local:finite-coherent-layer}
With $p\ge2^{32}$ and $0\le t\le12/5$ as fixed above, on the explicit gates $m>0$ and $c<1$, the normalized physical
layer operator of the law in \cref{lem:imp:coherent39-law},
represented on the fixed uniform six-distinct induced half and
written $T$, satisfies
\begin{equation}\label{eq:merge-local-I-37}
 -g_{21}rI\preceq T\preceq(1+B_{21}q)I,
\end{equation}
where
\begin{equation}\label{eq:merge-local-I-38}
 g_{21}=\frac{z(219/100)\Gamma+Cr}{m}
                  +\frac{C^2qr}{m^2},\qquad
 B_{21}=\frac{C^2q}{m^2(1-c)}.
\end{equation}
Its marginal defect (the difference between its actual half
marginal and the fixed chart marginal, as a density minus $1$)
has $L^2$ norm at most $\delta$, hence variance at most
$\delta^2$. Its pointwise marginal range ratio is at
most the old range ratio multiplied by $W/\eta$.
Evaluation of these positive majorants at the endpoint $p=X$
(the lower endpoint of the interval) is valid uniformly for every
prime $p\ge X$. The separate chamber inequalities and the
completion inequalities for the law itself (not a converted one)
remain to be checked.
\end{proposition}

\begin{proof}
We first justify every finite constant. A cyclic triangle has
density $d=8/(1-q^2)$ relative to its three-label cube and
exactly $(p+1)/4$ completions of a directed pair. The inherited
six-cycle expansion gives \eqref{eq:merge-local-I-32} with
\begin{equation}\label{eq:merge-local-I-39}
 \|E\|\le(6+51d)q<500q,\qquad
 \|rFM_0F^*\|\le\kappa(1+5r)r,\quad
 \kappa=(\sqrt3+\sqrt7)/2 .
\end{equation}
The twelve star entries are six evaluation/evaluation and six
evaluation/pair entries. The positive row and column bound
is $M$, and $0\le\Phi\le64$, so their squared norms are at
most $64M<600$; also $\|\Phi\|<2$.
Consequently the old two-component root equals its constant
plus leading star operator with remainder at most $1100q$,
since $1000+[\kappa(1+5r)+500r]^2<1100$.

For completeness the finite product bounds used next hold
also when a raw symmetric-difference support is the full
triangle. If $|S|=2$ and $1\le|T|\le3$, the normalized pair
star times $\psi_T$ has decomposition $h_{ST}e_0^*+R_{ST}$.
The exact diagonal image is
$(\mathbb E\psi_S\psi_T-m_Sm_T)/\sqrt{1-q}$.
Off the auxiliary diagonal, a matching/path or restricted
$K_{2,2}$ map costs at most $2\sqrt d\,r<6r$.
The mean-zero diagonal part is below $2r$, using
$D_H^*D_H=(1-2q)P_H+(1-q)J$ and the at-most-three removed
roots. The constant-input part is at most $2r$.
Thus $\|R_{ST}\|\le10r$, $|h_{ST}|\le2$, and
$S\ne T$ gives $|h_{ST}|\le3r$. A raw singleton gives a
whole product at most $13r$. A coordinate times a raw pair
or triple has norm at most $\sqrt d\,r<3r$, by fixing that
coordinate and retaining a free Paley transform.
Every whole product has norm at most two.
The coarser constants $100r,4r,2$ will suffice.

There are 24 one-component leading terms in the linear
rooted expansion. Each costs at most $4q$. Of the 144
two-component terms, 126 have a free evaluated leaf and cost
at most $10q$ each; fixing that coordinate leaves conditional
density at most $4/(1+q)<4$. The 18 exceptional
evaluation/pair against pair/evaluation terms have residual
at most $400r+10000q+8q^2<500r$ before their old factor $r$.
This includes the exact double-$e_0$ contraction. The 24
remainder/star placements cost at most $24\cdot4\cdot500q$;
the two remainder/rank terms cost $1000rq$, and the
remainder/remainder term costs $250000rq$. Their sum is
below $100000q$ per trace. Summing the absolute coefficient
mass $L$ proves the linear price used in \eqref{eq:merge-local-I-35}.

For the quadratic term, each of the 1521 valid crossed
factorizations has two raw two-label maps at most $6q$.
Its positive Gram CS bound is at most
$600\cdot9\cdot4q=21600q$. The diagonal two-sided raw-trace
bound is at most $6100q$: opposite components use two raw
transforms of norm at most $2r$, giving $2400q$, while a
common component uses its rank,
twelve star entries and $500q$ remainder, followed by
the other component of norm less than two.
The evaluation endpoint controls its exceptional diagonal
image. Holes cost at most $32L^2q$.
Collect the one-sided pair terms only after conditioning
both halves to be distinct; then they are exactly $-q$
and have absolute coefficient mass $8/25$.
Thus the full induced-half covariance error is at most
\begin{equation}\label{eq:merge-local-I-40}
 \frac qh\{(21600+6100+32)L^2+4(8/25)\}
 <200000q/h.
\end{equation}
Before conditioning a pair spanning the two triangles may
have trace $1_{x_i=x_j}-q$, so no unconditioned full-norm
claim is used here.

For collisions in the actual kernel, twelve equalities between
labels at different bases make the untilted (old) root kernel
vanish. The six remaining twin events
have positive total price $384Mq/h$; the eighteen
cross-component events cost $72d^2(1+q)^5q/h$.
Multiply their union by $W$. Partial root repetitions
still satisfy the even/odd degree bound in \eqref{eq:merge-local-I-30}.
Complete cancellation requires a three-matching between
the halves. There are at most
$\binom63^2\,3!=2400$ such matchings. The input-half cube
density is at most $d^2/h$, the root kernel at most 4096,
and each matching has mass $q^3$. The unrestricted score
bound $L/r$ therefore gives exactly $qC_{\rm exc}$.
Positive unions include all further collisions and
intersections; no coupled mask was removed inside a norm.

The finite actual-star Gram is at most $1+182r$.
Indeed every conditional vertex has at least $(p-11)/4$
choices, and the common actual-star identities give
off-diagonal Gram entries at most $7r$ among 27 active arms.
The fifteen mean arms have Gram at most
\begin{equation}\label{eq:merge-local-I-41}
 \frac{(1+5r)(1+10r+5q)}{1-9q}\le1+16r:
\end{equation}
grouping by singleton/pair support in one triangle, using its
raw six-arm frame, and then using the pointwise character Gram of
the other triangle. Both are bounded by $\Gamma$.

Each mean arm has expectation at most $100q/h$.
Before half conditioning its raw mean is zero when the support
meets the two triangles in $1$ and $2$ labels, and $q^2$ when it meets them in $2$ and $2$ labels. Removing nine equalities
costs $27qr$, except complete two-pair cancellation of
probability at most $18q^2$; division by $r h$ gives
$(27+19r)q/h<100q/h$. With mean synthesis norm at most two,
coefficient norm at most $\sqrt\sigma<1/2$, and fifteen
arms, centering the raw-mean contribution to $2tH$ costs
less than $10q$. Nonconstant auxiliary arms already kill
constants by translation invariance.

We have therefore proved the exact decomposition
$T_{\rm raw}=zJ+L_{\rm cen}+R$ with
\begin{equation}\label{eq:merge-local-I-42}
 L_{\rm cen}1=0,\quad
 \|R\|\le Cq,\quad
 L_{\rm cen}\succeq-z(219/100)\Gamma rI,\quad
 \|L_{\rm cen}\|\le4z\Gamma r .
\end{equation}
The lower order uses the 54-dimensional rational gate, the
separate invariant mean gate and \cref{imp:local:coherent-parameter}; the upper norm follows
from $\kappa+\lambda\sqrt\sigma<4$.
The actual raw scalar mass is at least $m$.
After normalization the constant entry is exactly one,
the off-diagonal vector has norm at most $\delta$, and
the centered block has norm at most $c$.
By \cref{lem:block-completion}, the positive constant square costs
at most $\delta^2$ below, and the separated upper blocks cost
$\delta^2/(1-c)$. These are precisely \eqref{eq:merge-local-I-38}.
The pointwise weight bounds $\eta\le w\le W$ prove the
range assertion without a pointwise closeness claim.

Finally every numerator after its indicated power of $q$
or $r$ is removed has nonnegative coefficients and
increases with $q,r$; $h,m,1-c$ decrease on their positive
gates. Rational outward endpoint evaluation is uniform.
The independent checker rebuilds the 39-support dictionary,
both collected diagonals and 136 scalar quantities at
eight endpoints, authenticating the separate matrix and
triangle-product certificates. Its parameter replay checks 72
endpoint/parameter cases and 1296 rational quantities, recovering
the values at the fixed parameter $t=t_*$. At $X=2^{45}$ the selected
$t=67/40$ gives $g_{21}<2.869198$ and $B_{21}<0.460808$, compared
with $g_{21}<3.163793$ and $B_{21}<0.947552$ at $t=t_*$.
\end{proof}

At the selected exact parameter $t=67/40$, the chamber strength is
$\lambda=402000/179557$ and a rational matrix gate improves the
convex chamber coefficient to $5.41674864$.
The checker \texttt{check\_coherent\_39\_selected\_matrix.py}
reconstructs the physical dictionary and directed reflected matrix,
proves all 56 rational LDL pivots positive, and verifies all 1596
exact reconstruction identities. The invariant mean order remains
below $3.2$. The parameter interval above is used for the finite
remainder and layer estimates; this stronger matrix gate is used
only at its stated exact parameter.
\subsection{A finite chamber price for the coherent law}
\label{imp:env:finite-coherent-section}

This is the corresponding numerical evaluation at the chamber
cut, for the same law \cref{lem:imp:coherent39-law} and with the
same convention: all norms are in $L^2$ of the fixed uniform
chamber half law, and no extra change-of-measure factor is
charged here. Throughout this subsection
\begin{gather*}
 q=p^{-1},\quad r=\sqrt q,\quad p\equiv3\pmod4,
 \quad p\ge2^{40},\\
 0\le t\le\frac{12}{5},\quad\eta=\frac1{100},\quad
 \sigma=\frac{13}{75},\quad s=z(t)=1+\eta+t^2\sigma.
\end{gather*}
The weight of the law (before normalization) is
$w=\eta+(1+tH)^2$. Recall that the $39$ coefficients in $H$ are
$\pm1/15$, of total absolute mass $\ell=13/5$ (denoted $L$ in the preceding subsection). The degree count
and Weil estimate in \eqref{eq:merge-local-I-30} give
\begin{equation}
 |H|\le B_H:=14+r
 \label{imp:env:coherent-score-bound}
\end{equation}
on the support of the uniform law.

Use the exact rectangle of the uniform law
(\cref{prop:l20-rectangular}) and its probability-normalized
decomposition
\[
 A=a\otimes1+r(\sqrt{30}Zv^*+2we^*)+E,
 \qquad \|E\|\le N_0,
\]
with $N_0$ the nonstar bound of
\cref{imp:env:finite-coherent-chamber}, and where the repeated
letter $w$ in this rectangle denotes its
four-coordinate feature function, not the weight $w$ above.
Write $a_+$ for a pointwise upper bound on $a$, $a_2$ for its
$L^2$ bound, and put
\[
 \theta=(1-q)(1-3q)(1-7q),\qquad
 h=(1-4q)(1-5q),\qquad D_6=\frac{16}{\theta h}.
\]
The row and column squared norms of $A$ are at most $256a_+$.
Let $K_\nu$ be its swapped product after compression to the old
six-distinct induced half.  A norm-bounded independent-half remainder
costs at most $1/h$; an unchanged leading feature has its original
coefficient after compression.

\begin{lemma}[Full coherent covariance]
\label{imp:env:coherent-full-covariance}
One has
\[
 \|K_\nu\odot H^2-\sigma K_\nu\|\le E_2/h,
\]
where
\begin{align}
 E_2=q\biggl[256a_+\biggl\{
       \frac{20992}{225\theta}+\frac{2108}{225(1-q)}\biggr\}
   +\biggl(\frac{32}{225}+\frac8{25}+\frac{2028}{25}\biggr)a_2^2
      \biggr]<26400q.\label{imp:env:coherent-E2}
\end{align}
This is a full operator estimate, including its constant sector.
\end{lemma}

\begin{proof}
For each ordered pair of score supports, keep their common auxiliary
labels and split the diagonal from its complement.  On the complement,
use the cheaper of the two valid whole-rectangle Cauchy--Schwarz
factorizations.  A different-support core product costs
$4q/\sqrt\theta$ and an equal-support product twice that; the
private costs are $q/\sqrt{1-q}$ and twice that.  The exact ordered counts of support pairs, by
core and private product coefficient, are
\[
 \begin{array}{c|rrr}
 \text{core product coefficient}&16&32&64\\
 \text{number}&804&56&41\\ \hline
 \text{private product coefficient}&1&2&4\\
 \text{number}&184&240&196.
 \end{array}
\]
Multiplying by the rectangle row or column square and by $1/225$
gives
\[
 256a_+q\left\{\frac{17280}{225\theta}
                  +\frac{1448}{225(1-q)}\right\}.
\]
The operator Cauchy--Schwarz uses the two output Grams for the
swapped rectangle, or the two input Grams in its private version.
It does not introduce a second copy of either auxiliary label.

On the auxiliary diagonal, collect equal symmetric differences before
taking absolute values.  Every nonzero one-sided difference has size
two.  Its trace in the induced half is exactly $-q$; the absolute
coefficient mass is $4/25$ on each side.  For the two-sided differences,
the private raw-star maps give weighted cost $44/15$, and the remaining
core maps give $3712/225$.  A core singleton has norm $r$; a higher
raw core star has norm at most $4r/\sqrt\theta$.  The only unmatched
terms are sixteen core-singleton/private-singleton terms of total
mass $32/225$.  Their raw trace is $1_{x=y}-q$, and the accepting
rectangle vanishes at $x=y$, so their operator is exactly $-qK_\nu$.
Finally the repeated-root holes cost at most
$12\ell^2q\|K_\nu\|$.
These are precisely the remaining terms in
\eqref{imp:env:coherent-E2}.
The one-sided identities are asserted after induced-half conditioning;
before conditioning a core/private pair may coincide.  The other
remainder estimates are compressed with the single factor $1/h$.
\end{proof}

\begin{lemma}[Finite coherent linear error]
\label{imp:env:coherent-linear-finite}
With $E_1=12500q$,
\[
 \|Q(K_\nu\odot H-H)Q\|\le E_1/h,
 \qquad \|K_\nu\odot H\|\le(12r+E_1)/h.
\]
\end{lemma}

\begin{proof}
Here are explicit rounded product bounds sufficient for the assertion.
The inherited positive endpoint majorants give
\[
 N\le282q,\qquad \|a-1\|_2\le4.5r,
 \qquad a_+,a_2,\|Z\|<1.001,
 \qquad \|K_\nu-J\|<15r.
\]
For higher core supports $S,T$, split the normalized product map by
the constant diagonal vector
$e_0=(p1_{u=v}-1)/\sqrt{p-1}$, its mean-zero diagonal, and its
off-diagonal complement.  Set
\begin{gather*}
 b_1=0,\quad b_2=q,\quad b_3=2r,\quad b_4=2r+q,\\
 D_2=\frac2{1-q},\quad
 D_3=\frac4{(1-q)(1-3q)},\quad D_4=\frac{16}{\theta}.
\end{gather*}
For $S\ne T$, writing $c=|S\cap T|$ and $d=|S\triangle T|$,
the three pieces are bounded respectively by
\[
 \sqrt{D_{|S\cup T|}},\qquad
 \frac{1+cr+b_{|S|}+b_{|T|}}{\sqrt{1-2q}},\qquad
 \frac{b_d+cq+b_{|S|}b_{|T|}}{r\sqrt{1-q}}.
\]
For $S=T$ double the first bound, omit the last, and remove the
constant image $1/[r\sqrt{1-q}]$.  Its remaining variation is at
most $(|T|q+b_{|T|}^2)/[r\sqrt{1-q}]$.  These bounds follow from
the two-Paley matching/path extraction or the restricted $K_{2,2}$
bound, and the exact Gram of the projected auxiliary diagonal.
Their orthogonal-input square sum, the eleven-star expression for
$Z$, and its singleton correction give residual bound 22 and
exceptional coefficient at most $.259/r$.  The four-singleton aggregate
has product bound 9.  Each expression is a positive endpoint majorant;
the displayed rational constants hold already at $p=2^{40}$.

There are thirteen pure-core terms (whose supports at this cut lie in the core labels of the halves).  For each, the rectangle expansion
prices the three rank and star classes, the three classes containing
a nonstar remainder, and the raw invariant means by the following
coefficients of $q$:
\[
 9.1,\quad278,\quad400,\quad565,\quad801,\quad1,\quad61.
\]
For clarity the nontrivial rounded expressions are
\begin{align*}
 278&>2(1+4.5r)(22\sqrt{30}+18),\\
 400&>30r\,22^2+60(.259)22
       +\frac{30(.259)^2r}{1-q}
       +4\sqrt{30}r\,22\cdot9
       +4\sqrt{30}(.259)9+4r\,9^2,\\
 801&>2\cdot282\{\sqrt{30}(.259+22r)+18r\}.
\end{align*}
The factor $1+4.5r$ is the invariant weighted-star correction from
$a-1$.  The large constant output of a pure-core rank/star term is
killed by $Q$ on that half; its nonconstant variation was included
in the residual bound 22.  The retained uncentered rank/star norm, after
weighting all thirteen terms, is below $6r$.  The mean contribution
uses $(2+r)^2\|K_\nu-J\|r<61q$.  The sum of these seven prices is
below 2300.

The other twenty-six terms are mixed.  Weighted raw-star synthesis
gives a rank/rank cost of at most $91q$.
The free core and private transforms give a rank/star cost of $61q$.
For the two cross-arm star/star
terms, the $w,v$ triple map has a free Paley leaf and norm below $6r$,
while $Z,e$ has norm below 1.001; their price is 132.  Equal $Z,v$
arms have exceptional amplitude at most $.259$ and residual $24r$;
equal $w,e$ arms have amplitude at most $.501$ and residual $6r$.
The double exceptional contractions are exactly $qQ/(p-1)$, not
products of their separate norms.  These prices are 375 and 26.
The nonstar/star chain applies the input product, then $E$ fiberwise,
then the output product, keeping both shared labels.  Their norms
are at most 1.04 and 1.001, so both placements cost less than
\[
 2(\sqrt{30}+2)(1.04)(1.001)282q<4500q.
\]
Nonstar/rank and nonstar/nonstar cost 566 and 1.  Thus the mixed
price is below 6000.  The total is bounded by
$(13\cdot2300+26\cdot6000)q/15<12500q$.

For the full-norm assertion, the raw score coefficient matrix has
absolute row sum at most $1/3$.  Cube star ranges are orthogonal,
so compression gives $\|H\|\le D_6r/3<5.334r$.
Adding the retained rank/star norm below $6r$ and the error just
proved gives the stated $12r+E_1$ bound, with its compression price.
\end{proof}

We state the final finite scalar interface explicitly.  Let $N_0,D,
\epsilon,\chi$ be the authenticated old raw nonstar, positive collision,
scalar-normalizer and normalized marginal-variance bounds:
$N_0=N_{\rm ns}$, $D=D_{\rm bal}$, $\epsilon=\eps_Z$ and
$\chi=\chi_p^+$ in the notation of
\cref{prop:l20-completion,prop:l20-collisions}. Put
\begin{align}
 P_*&=144\cdot65536D_6q^3,\nonumber\\
 C_w&=[\eta+(1+tB_H)^2]D
       +P_*[\eta+1+2t\ell/r+t^2\ell^2/q],\nonumber\\
 L&=\frac{(12r+E_1)/h+B_HD+\ell P_*/r}{1-\epsilon},\nonumber\\
 V&=\frac{E_2/h+(B_H^2+\sigma)D+(\ell^2/q+\sigma)P_*}
                  {1-\epsilon},\nonumber\\
 z&=s-2tL-t^2V,\qquad
 \delta=\frac{s\sqrt\chi+2tL+t^2V}{z}.
 \label{imp:env:coherent-normalization}
\end{align}
(From here $z$ denotes the proved lower bound $s-2tL-t^2V$ for the scalar mass whose leading value was called $z(t)$ above; the letter $L$ in $z$ is the linear error term of the display, not the block budget.) Let $\Gamma=1+434r$. Let $g_*$ be any number such that both the
nontrivial and invariant branches at the specified $t$ have
leading negative coefficient at most $g_*$; the value
$g_{\rm ch}(t)$ of \cref{imp:env:coherent-parameter} may be
used, or a sharper certified coefficient
at the specified $t$. In particular the selected
$t=67/40$ uses $g_*=5.41674864$. Define
\begin{align}
 E_*&=sN_0+2tE_1/h+t^2E_2/h+C_w,\nonumber\\
 G&=\frac{s g_*\Gamma r+E_*}{(1-\epsilon)z},\qquad
 c=\frac{s\sqrt{30}\Gamma r+2tD_6r/3+E_*}{(1-\epsilon)z}.
 \label{imp:env:coherent-final-prices}
\end{align}

\begin{proposition}[Finite chamber-cut bound and marginal data for the coherent law]
\label{imp:env:finite-coherent-chamber}
The explicit gates $z>0$, $c<1$, and $\delta<1$ hold on the stated
domain.  The normalized physical operator of
\cref{lem:imp:coherent39-law}, on the fixed old induced half and written $T$, satisfies
\[
 -(G+\delta^2)I\preceq T
 \preceq\left(1+\frac{\delta^2}{1-c}\right)I.
\]
Its marginal variance is at most $\delta^2$, and its marginal range
ratio is at most the old ratio times
$[\eta+(1+tB_H)^2]/\eta$.  Every endpoint price is uniform for all
primes $p\equiv3\pmod4$ above the endpoint.
\end{proposition}

\begin{proof}
Complete character-root cancellation needs three cross identifications
for each of the twelve degree-six traces, or four for each of the
three degree-eight traces.  There are at most
$12\cdot3!+3\cdot4!=144$ matchings.  The raw root is bounded by 65536
and the input half has cube density at most $D_6$, giving the positive
exceptional row/column mass $P_*$.  Outside it, cancellation only
reduces the Weil price in \eqref{imp:env:coherent-score-bound}; everywhere
$|H|\le\ell/r$.  Thus $C_w$ bounds the difference between the soft kernel and
the kernel with the distinctness conditions imposed, including all intersections.  The same argument and the two lemmas give the full
normalized-old prices $L$ and $V$.  The new scalar mass is at least $z$,
and its centered row has norm at most $\delta$.

The general actual-star Gram bound applies to all 63 nonempty supports.
Every distinct-feature entry is at most $7r$, giving $\Gamma$.
The exact matrix inequality of
\cref{imp:env:coherent-parameter} (its certified gate) gives
the nontrivial negative leading price $s g_*\Gamma r$. In the
invariant sector the bound from the core mean features
price is $16/(\theta h)$, and the mixed price is at most
\[
 \frac{[1+q/(1-q)+2qr/(1-q)](1+45r+60q)}h.
\]
The latter follows by enlarging the core character list to all 15
nonempty supports: each distinct-support entry is at most $3r+4q$,
so its row sum is at most $1+42r+56q$.  The respective negative
coefficient orders $-7/100$ and $-21/100$ keep this invariant branch
strictly below the nontrivial branch.  The full positive centered
price follows from $\|M_{\rm old}\|=\sqrt{30}$ and
$\|H\|\le D_6r/3$.  This proves the centered bounds $-G$ and $c$.
The constant-block bounds of \cref{lem:block-completion} cost at most
$\delta^2$ below and $\delta^2/(1-c)$ above.  The pointwise multiplier lies between
$\eta$ and $\eta+(1+tB_H)^2$, proving the range statement.

After substituting $q=r^2$, all new normalized numerator majorants
are positive increasing functions of $r$, while $z$ and $1-c$ decrease.
For example the three exceptional terms in $C_w/r$ have powers
$r^5,r^4,r^3$.  The inherited majorants have the same endpoint
property.  Thus the gates and orders extend uniformly above the
endpoint.  Exact coefficient collection, finite product gates and
an independently assembled scalar replay accompany the proof (\cref{app:certificates}).
For the endpoint parameter $t=t_*$ at $2^{45}$, with $g_{20}$
the negative coefficient in $T\succeq-g_{20}rI$ and $B_{20}$ the
excess in $T\preceq(1+B_{20}q)I$ (both formed from $G,\delta,c$
by the formulas above, as $g_{21},B_{21}$ are for the layer cut),
the values are $g_{20}<5.441282211737$ and
$B_{20}<1100.871034665$.
These are local interfaces; the layer price, native completion and
full scalar comparison are checked separately.
\end{proof}

\begin{proposition}[The coherent score on a single collision]
\label{imp:env:coherent-single-pin}
For the same $39$-trace coherent law, put
\[
 B_c=\frac{64}5+\frac85r,\qquad
 B_s=10+\frac{11}5r,\qquad
 B_x=\frac{26}3+\frac{31}{15}r,\qquad B=14+r.
\]
Let $A_c,A_p,A_m$ be the positive collision bounds
\[
\begin{split}
 e_2&=r+48q,\quad e_3=5r+96q,\quad e_4=17r+192q,\\
 \Lambda_c&=(1-r-24q)^4(1-6q)^2,\\
 A_c&=\frac{32q(1+e_3)^2(1+e_4)^3}{\Lambda_c},\\
 A_p&=\frac{32q(1+e_2)(1+e_3)^2(1+e_4)^2}{\Lambda_c},\\
 A_m&=\frac{(1+e_4)^6}{\Lambda_c}
 (3712q^2+81920q^3+1245184q^4+12582912q^5+67108864q^6).
\end{split}
\]
With $w(x)=\eta+(1+tx)^2$, replace only $C_w$ in the negative
price of Proposition~\ref{imp:env:finite-coherent-chamber} by
\[
 C_{\rm pin}=A_cw(B_c)+A_p[w(B_s)+w(B_x)]+A_mw(B)
 +P_*[\eta+1+2t\ell/r+t^2\ell^2/q].
\]
Retain the previous $z,\delta,c$ and all positive, variance and range
prices. The resulting lower coefficient is uniformly valid for
$p\ge2^{40}$ and each fixed parameter $t$ for which the leading matrix bound has been certified.
\end{proposition}

\begin{proof}
At the chosen chamber cut, numbering the label $(j,\ell)$ by $2j+\ell$, the coordinate sets are
\[
 \text{cores: }\{0,1,4,5\},\ \{6,7,10,11\};
 \qquad \text{private pairs: }\{2,3\},\ \{8,9\}.
\]
Only their sixteen core/core and four
private/private coincidences survive the soft accepting kernel.
For exactly one equality $x_a=x_b=u$, use
$\chi(v-u)^2=1-1_{\{v=u\}}$ in every score monomial and collect the
remaining supports.  In the score divided by $r$, the deleted
evaluation contributes at most $r$ times its absolute coefficient.
No remaining support has fewer than four labels.  For a remaining support of degree $d$, the even-degree
character price is $d-2+r$ and the odd-degree price is $d-1$.
Exact collection gives eight pins with bound $B_c$, six with $B_s$,
and six with $B_x$.

The eight $B_c$ pins join the output $a$ pair to the input $a$ pair,
or output $c$ to input $c$.  The other two classes are
\[
\begin{array}{c|l}
 B_s&(0,6),(1,7),(4,10),(5,11),(2,8),(3,9)\\
 B_x&(0,7),(1,6),(4,11),(5,10),(2,9),(3,8).
\end{array}
\]
For $p\ge2^{40}$, $B_c$ dominates the other two prices, so the original
mutually exclusive core orientations retain their total incident
charge $32q$. For each output private coordinate, the two input
private pins have the same incident cyclic-extension weight, at
most $16$, and one belongs to each of the $B_s,B_x$ classes.
Their row charges are therefore $32q$ each. The counts of
extensions of a pinned pair to full labellings (the extension
histograms) are unchanged within the two groups and give
$A_c,A_p$; transposition
gives the same column bounds.

The strata with exactly one coincidence are disjoint from the set with at least two coincidences, already bounded by $A_m$.
Keep the original $A_m$ with the generic score bound $B$, and the entire
exceptional mass $P_*$ with the worst score $\ell/r$.
Schur's test gives the displayed collision price.
Recompute the numerator of $G$ in \eqref{imp:env:coherent-final-prices}, with $C_{\rm pin}$ in place of $C_w$ in $E_*$, using only positive terms, retaining
the old denominator and constant-block charge.  This proves the
same-law substitution without subtracting monotone majorants.
All normalized numerator prices increase with $q$ and $r$, and the previous
denominators decrease on their gates, proving endpoint uniformity.
The independent replay verifies the 20 support collections,
180 finite identities and all 35 endpoint/parameter cases.
At $p=2^{45}$ and the rational parameter $t=3337/2048$ (a
rational below the endpoint $t_*=12/5$ at which this refined
check is run; the selected construction uses $t=67/40$ with the
coefficient above), the refined coefficient is
$g_{20}<5.434524687965$, with the positive and marginal prices unchanged.
\end{proof}


\subsection{Values of the bounds at the seven interval starts}\label{subsec:positive-root-caps}

The preceding subsections proved bounds for the tilted law in
terms of quantities evaluated at the start of each interval. The
following certificate records those values at the seven starts,
checked by exact and outward-interval arithmetic; the table at its
end is the form in which later sections use them.

\begin{certificate}[The positive roots]\label{cert:positive-root-caps}
The chamber caps for the selected individual-square (pair-score) laws combine
\cref{imp:local:collision-gap,imp:env:finite-gap-chamber,imp:env:physical-interpolation};
their layer prices use \cref{imp:env:layer-resummation} where it improves the
retained finite price. The last band uses
\cref{imp:local:coherent-parameter,imp:env:coherent-parameter,imp:env:coherent-single-pin},
with the exact matrix gate at $\ulaw=67/40$.
All prices concern the same tilted law at both cuts, and that law's
actual marginal enters the restoration. Every positivity or
contractivity condition required above --- in particular
$z_->0$, $\Delta_g>0$, $\rho_{\rm in}>0$, $c_{\rm ch}<1$,
$d_\ulaw<1$ and $gX^{-1/2}<1$ at each endpoint $X$ --- is checked
in outward interval arithmetic, with the uniform endpoint
majorants proved above.

The exact selected data are the \texttt{local.spectra} and
\texttt{local.further\_root\_law} entries of
\path{research/push_2026_10_25/global/refined_parent.json}. They agree
exactly with the native data authenticated by
\path{research/improvements_2026_10_25/verify.py}; only scalar charges
are reanchored by \path{research/push_2026_10_25/verify.py}.
The former reruns independent collision, pair-gap, interpolation,
resummation, coherent covariance, selected matrix, finite chamber,
finite layer and single-pin checkers. Its full tier regenerates the
native completion. The latter checks exact equality of all native
records and reruns the independent closing and restoration assemblies.
At $L_{\rm loc}=585$, the exported prices, rounded upward, are:
\begin{center}\scriptsize
\begin{tabular}{@{}llllll@{}}
\toprule
piece start & law, $\ulaw$ & cut & $g$ & $c^+$ & block cap\\
\midrule
$2^{30}$ & individual squares, $711/4096$ & L20 & $5.476619678296$ & $25.456659470084$ & $5.476623331076$\\
 &  & L21 & $6.313078885225$ & $3.988889092310$ & $6.313079545011$\\
$2^{31}$ & individual squares, $1021/4096$ & L20 & $5.474335244881$ & $27.623843030000$ & $5.474337162024$\\
 &  & L21 & $6.315256713466$ & $3.962812801799$ & $6.315257030740$\\
$2^{32}$ & individual squares, $1455/4096$ & L20 & $5.471257426364$ & $30.397461416086$ & $5.471258447645$\\
 &  & L21 & $6.312562621834$ & $3.927444212402$ & $6.312562774077$\\
$2^{34}$ & individual squares, $833/1024$ & L20 & $5.464133834236$ & $36.835698528597$ & $5.464134125055$\\
 &  & L21 & $6.307298870667$ & $3.877354549660$ & $6.307298906003$\\
$2^{36}$ & individual squares, $3851/4096$ & L20 & $5.460775686025$ & $37.085210712622$ & $5.460775755114$\\
 &  & L21 & $5.426082230198$ & $1.143781717023$ & $5.426082232316$\\
$2^{40}$ & individual squares, $4049/4096$ & L20 & $5.457521105807$ & $37.019255725148$ & $5.457521109684$\\
 &  & L21 & $4.646587567237$ & $0.075183415124$ & $4.646587567244$\\
$2^{45}$ & coherent 39, $67/40$ & L20 & $5.434450251962$ & $983.324753609660$ & $5.434450254811$\\
 &  & L21 & $2.869197519095$ & $0.460807795864$ & $2.869197519096$\\
\bottomrule
\end{tabular}
\end{center}
The seven per-prime marginal variances and range ratios are recorded as
exact rationals in the corresponding rows of the data file named above and
displayed in
\cref{tab:data:positive-roots}. They are bounds in the fixed uniform
chamber half, with its actual root marginal retained. No return to a
uniform root is made in computing the restorations.
\end{certificate}

The finite tables retain the larger choice
$\mathfrak b_{\rm cert}=6+51d_\mu$ in the layer formulas and the
corresponding certified remainder prices. Since
$\mathfrak b\le\mathfrak b_{\rm cert}$, these remain valid
majorants for the same laws. The sharper triangle estimate and
\cref{lem:native-pair-linear} do not change the chosen weights or
permit an uncharged change of the actual half marginals.

The block caps of L20, rounded upward, give $w_{\rm D}=w_{\rm S}$,
and those of L21 give $v_{\rm L}$ (\cref{rem:middle-weights}). The
restorations use actual budget $570$ and the respective maximum block
cardinalities $\lfloor L/\log X\rfloor$ ($L=570$), namely
$27,26,25,24,22,20,18$ at the seven starts; local operator tensorization
uses budget $585$.

\begin{remark}[The family of tilts]\label{rem:positive-root-family}
The root $g(\lambda)$ of \eqref{eq:root-cubic} is the limit of the chamber coefficient
$g^\ulaw_{20}$ as $p\to\infty$ (\cref{prop:root-chamber}: $\Gamma_{20}\to1$, $z_{\rm raw}\to s$).
It decreases from $g(0)=\sqrt{30}$, the limit of the uniform root, to about $5.466411$ at
$\lambda=1/2$, $5.456962$ at $\lambda=24/25$ ($\ulaw=3/4$), $5.456161$ at $\lambda=1$ ($\ulaw=1$,
the largest value of $2\ulaw/(1+\ulaw^2)$).
For comparison, the leading matrix of a square-of-the-average weight
has strength $6\ulaw/(3+\ulaw^2)$: the same cubic gives $5.443543$ at
$\lambda=4080/2467$ ($\ulaw=51/40$) and $5.442090$ at its maximum
$\lambda=\sqrt3$. These matrix values explain the historical interest in
that family; a finite law must also pay its mixed-product errors. The gain is at order one; the costs are at order $r$: the layer coefficient
tends to $\sqrt3+\sqrt7$ for every $\ulaw$, the half marginal moves, and the range ratio
grows like $(1+\ulaw(2+r))^2/(1-\ulaw c_{\rm cond})^2$ and enters the restorations through the
per-block product of the per-prime range ratios, bounded by
$M_{\Blk}\le R_\ulaw(X)^n$ in \cref{lem:record:root-marginal}(iv). The selected seven middle pieces and their parameters are those of
\cref{def:positive-root}. The coherent last-band law has its own full
chamber and layer estimates above; the one-layer root is unchanged.
\end{remark}

\paragraph{Consequences and what to read next.}
\Cref{thm:weil} (the scores); \cref{lem:layer12-geometry} and \cref{lem:marginals} (the
symmetries); \cref{lem:l20-star-matrix,prop:l20-completion,prop:l20-rank4,prop:l20-rectangular,prop:l20-collisions}
(the star matrix and the frame cap, the normalizer, the rank-four Gram, the rectangle and the
collision cap of the uniform root);
\cref{lem:general-star,cor:star-matrix,rem:realified-gates} and \cref{lem:spec:affpair}(d)
(stars); \cref{cor:raw-pair-maps,lem:swapped-rectangle,lem:l20-star-mean} (pair maps, swapped
products, star means), with \cref{thm:dl:K}, \cref{lem:matching:kernel} and
\cref{lem:paley-triple} (witness kernels and completion counts);
\cref{lem:record:weighted-moment} (the fluctuation of the extension count);
\cref{lem:restriction} (the actual half);
\cref{lem:tilt-marginals} and \cref{prop:matching:positive}(a),(a$'$) (mass, marginal, range and
the constant block); \cref{lem:record:root-marginal}(i),(ii) and
\eqref{eq:record:root-count} ($R_p$, $\chi^+_p$); \cref{lem:spec:cutkernel}(f) and
\cref{lem:spec:tensorcomp}(b) (cut operators of a reweighted law, block caps);
\cref{lem:majorant} (evaluation at an endpoint); \cref{lem:record:local-domain} (the block
budget). Certificate:
\cref{cert:positive-root-caps}.

\paragraph{Hypotheses verified in this section (\cref{tab:hypotheses}).} the seven positive-root bands
supply the root inputs of ``(CP1) middle'': positivity, support, native
sorting symmetries, uniform single-slot marginals, and finite chamber
and layer caps. Pair laws use \cref{lem:positive-root-properties};
the coherent law uses the physical dictionary, positive floor and
square-affine invariance of \cref{lem:imp:coherent39-law}. Their common finite interfaces
and actual marginal data are certified in \cref{cert:positive-root-caps}.
The completion is in \cref{sec:middle-frame}; restorations use
\cref{lem:record:root-marginal}(iv).
